%%%%%%%%%%%%%%%%%%%%%%%%%%%%%%%%%%%%%%%%%
%
% PSI Chair Thesis Template
% Version 20200913
%
% based on MastersDoctoralThesis.cls
% Version 2.5 (27/8/17)
%
% which was obtained from:
% http://www.LaTeXTemplates.com
%
% Version 2.x major modifications by:
% Vel (vel@latextemplates.com)
%
% This template is based on a template by:
% Steve Gunn (http://users.ecs.soton.ac.uk/srg/softwaretools/document/templates/)
% Sunil Patel (http://www.sunilpatel.co.uk/thesis-template/)
%
% License of this guide and the template
% CC BY-SA 4.0 (http://creativecommons.org/licenses/by-sa/4.0/)
%
% Exception 1: Some excerpts, figures, and tables that have been taken
% from the literature (denoted with a citation in the caption) are not
% covered by the above license. Permission to re-use and distribute
% these excerpts, figures, and tables must be obtained from the
% respective holder of the copyrights.
%
% Exception 2: Chapter 1 and Appendix C are based on content from the
% MastersDoctoralThesis template mentioned above, which is licensed under
% CC BY-SA 3.0 (http://creativecommons.org/licenses/by-nc-sa/3.0/)
%
% License of the PSIThesis.cls class file:
% LPPL v1.3c (http://www.latex-project.org/lppl)
%
%%%%%%%%%%%%%%%%%%%%%%%%%%%%%%%%%%%%%%%%%

%----------------------------------------------------------------------------------------
% PACKAGES AND OTHER DOCUMENT CONFIGURATIONS
%----------------------------------------------------------------------------------------

\PassOptionsToPackage{english,ngerman}{babel}
\documentclass[
11pt, % The default document font size is 11 (recommended), options: 10pt, 11pt, 12pt
oneside, % Two-side layout is recommended; uncomment to switch to one-sided
english, % replace with ngerman for German; not fully supported so far -- requires changes elsewhere
singlespacing, % Single line spacing (recommended), alternatives: onehalfspacing or doublespacing
% draft, % Uncomment to enable draft mode (no pictures, no links, overfull hboxes indicated)
%nolistspacing, % If the document is onehalfspacing or doublespacing, uncomment this to set spacing in lists to single
%liststotoc, % Uncomment to add list of figures/tables/etc to table of contents (not recommended)
%toctotoc, % Uncomment to add the main table of contents to the table of contents (not recommended)
parskip, % add space between paragraphs (recommended)
nohyperref, % do not load the hyperref package (is loaded in setup.tex)
%headsepline, % print a horizontal line under the page header
consistentlayout, % layout of declaration, abstract and acknowledgements pages matches the default layout
%final, % Uncomment to hide all todo notes
]{PSIThesis} % The class file specifying the document structure

% version of the guide
\def\tversion{v20200913}

% long-term stable URL to the thesis guide
\def\doiurl{https://doi.org/10.20378/irb-48428}
\def\githuburl{https://github.com/UBA-PSI/psi-thesis-guide}

%%%%%%%%%%%%%%%%%%%%%%%%%%%%%%%%%%%%%%%%%%%%%%%%%%%
%
% File: setup.tex
%
% This file is part of the PSIThesis.cls
% LaTeX documentclass
%
% The code in this file is made available
% under the following license:
%
% LPPL v1.3c (http://www.latex-project.org/lppl)
%
%%%%%%%%%%%%%%%%%%%%%%%%%%%%%%%%%%%%%%%%%%%%%%%%%%%


%----------------------------------------------------------------------------------------
%   UNIVERSITY OF BAMBERG COLORS
%----------------------------------------------------------------------------------------

% http://www.brandwares.com/RGBTintCalculator.php
% Base Color value obtained from UB Corporate Identity Manual
\definecolor{ubblue}{HTML}{00457D}
\definecolor{ubblue80}{HTML}{336A97}
% \definecolor{ubblue60}{HTML}{668FB1}
% \definecolor{ubblue40}{HTML}{99B5CB}
% \definecolor{ubblue20}{HTML}{CCDAE5}

\definecolor{ubyellow}{HTML}{FFD300}
\definecolor{ubyellow25}{HTML}{FFF4BF}

\definecolor{ubred}{HTML}{e6444F}

\definecolor{ubgreen}{HTML}{97BF0D}

\definecolor{gray75}{gray}{0.75}
\definecolor{gray50}{gray}{0.50}

% https://imagecolorpicker.com/color-code/a11e3b
% Base Color value obtained from Universitaet Tuebingen Logo
\definecolor{utred}{HTML}{A11E3B}
\definecolor{utred80}{HTML}{81182F}
\definecolor{utred60}{HTML}{611223}
\definecolor{utred40}{HTML}{400C18}
\definecolor{utred20}{HTML}{20060c}
\definecolor{triadicblue}{HTML}{1E3BA1}

%----------------------------------------------------------------------------------------
%   FONT SETUP
%----------------------------------------------------------------------------------------

%\usepackage[T1]{fontenc} % Output font encoding for international characters - do not use T1 encoding with luatex!
\usepackage[utf8]{luainputenc} % makes unicode characters like –, €, and ß work properly

% amssymb must be loaded before newtxmath to avoid this error:
% Command `\Bbbk' already defined
\usepackage{amssymb}
\usepackage[cochineal]{newtxmath} % must be loaded before fontspec

\usepackage[no-math]{fontspec} % allows us to use OTF/TTF fonts, but do not interfere with math (because we use newtxmath, which does support Cochineal)

% If you cannot use the cochineal font, uncomment the following lines to select
% the Crimson font. Note, however, that you'll have to take care of the math font
% on your own.
%
%\setmainfont[
% Path           = fonts/,
%    BoldFont       = {Crimson-Semibold.otf},
%    ItalicFont     = {Crimson-Italic.otf},
%    BoldItalicFont = {Crimson-BoldItalic.otf}
%]{Crimson-Roman.otf}

\setmainfont{Cochineal}[
  Numbers={Proportional,OldStyle},
  Style=Swash % for nice swashed Q letter, see https://golatex.de/spezeielle-opentype-features-in-fontspec-aktivieren-t19831.html
]
%\setmainlanguage{english}
\DeclareSymbolFont{operators}{\encodingdefault}{\familydefault}{m}{n} %  render numbers in cochineal, cf. https://tex.stackexchange.com/questions/398895/using-two-different-math-fonts-with-lualatex

% imitate the behavior of the cochineal package as follows:
% cf. https://tex.stackexchange.com/questions/448895/fontenc-vs-fontspec-with-xelatex
\DeclareRobustCommand{\lfstyle}{\addfontfeatures{Numbers=Lining}}
\DeclareTextFontCommand{\textlf}{\lfstyle}
\DeclareRobustCommand{\tlfstyle}{\addfontfeatures{Numbers={Tabular,Lining}}}
\DeclareTextFontCommand{\texttlf}{\tlfstyle}

% Exception: tables should use "lining figures" (all digits having same width)
\AtBeginEnvironment{tabular}{%
  \tlfstyle
}
\AtBeginEnvironment{tabularx}{%
  \tlfstyle
}


% monospace font, will be used in verbatim and listing environments
\setmonofont[
  Path           = fonts/,
    BoldFont       = {iosevka-ss04-bold.ttf},
    ItalicFont       = {iosevka-ss04-italic.ttf},
    BoldItalicFont       = {iosevka-ss04-bolditalic.ttf},
    Scale = 0.83 % manually determined value;
]{iosevka-ss04-regular.ttf}


% sans-serif font, will be used in the margins
\setsansfont[
  Path            = fonts/,
  BoldFont    = Roboto-Bold.otf,
  ItalicFont    = Roboto-Italic.otf,
  BoldItalicFont  = Roboto-BoldItalic.otf,
  Scale = 0.83 % manually determined value;
]{Roboto-Regular.otf}

\renewcommand{\familydefault}{\rmdefault}
\defaultfontfeatures{Ligatures=TeX}


%----------------------------------------------------------------------------------------
%   HEADINGS SETUP (CHAPTERS, SECTIONS, …)
%----------------------------------------------------------------------------------------

\usepackage[explicit]{titlesec}
\newcommand{\hsp}{\hspace{20pt}}

\setcounter{secnumdepth}{3}

% We use lining figures for headers (tlfstyle) because they fit better with uppercase letters than old-style figures.

% chapters have a vertical line between number and title
\titleformat{\chapter}[hang]{\Huge\bfseries\tlfstyle}{\color{black}\thechapter}{20pt}{\begin{tabular}[t]{@{\color{utred}\vrule width 2pt\hsp}p{0.85\textwidth}}\raggedright #1\end{tabular}}

% sections
\titleformat{\section}[hang]{\bfseries\large\tlfstyle}{{\color{utred60}\thechapter.\arabic{section}}}{1ex}{{\color{utred60} #1}}{}

% subsections
\titleformat{\subsection}[hang]{\bfseries\large\tlfstyle}{{\color{utred80}\thechapter.\arabic{section}.\arabic{subsection}}}{1ex}{{\color{utred80} #1}}{}

% subsubsections
\titleformat{\subsubsection}[hang]{\bfseries\tlfstyle}{{\color{utred}\thechapter.\arabic{section}.\arabic{subsection}.\arabic{subsubsection}}}{1ex}{{\color{utred} #1}}{}

% vertical spacing for headings ==============
\titlespacing*{\section}
{0pt}{7ex}{3ex}

\titlespacing*{\subsection}
{0pt}{4ex}{2ex}

\titlespacing*{\subsubsection}
{0pt}{4ex}{2ex}
% end of vertical spacing ====================

%----------------------------------------------------------------------------------------
%   TABLE OF CONTENTS SETUP
%----------------------------------------------------------------------------------------

% solution inspired from https://tex.stackexchange.com/questions/178510/how-can-i-reproduce-this-beautiful-table-of-contents
\usepackage{etoc}
\etocsetlevel{section}{2}
\etocsetlevel{subsection}{3}

\etocsettocdepth{section} % set to subsection for adding subsections to toc (not recommended)

\newlength{\tocleft}
\setlength{\tocleft}{2.5cm} % must be set to fit the innermargin defined in geometry (change only if you have changed the margins)

\newlength{\tocsep}
\setlength{\tocsep}{2em}

\usepackage{textcase}


\etocsetstyle{chapter}
   {}
   {}
   {\etocifnumbered
     {\makebox[0pt][r]
       % we use \etocthenumber instead of \etocnumber to avoid the href, which is part of \etocthenumber, messing with MakeTextLowercase
       {\textsc{\MakeTextLowercase\chaptername\ \MakeTextLowercase\etocthenumber}\hspace{\tocsep}}%
       \textbf{\etocname\kern1em\relax\etocpage}%
    }%
    {\textbf{\etocname\kern1em\relax\etocpage}}%
    \par\vspace{3ex}%
   }%
   {}

\etocsetstyle{section}
   {\vspace{-2ex}} % Muss von den 3ex aus Chapter abgezogen werden
   {}
   % see the comment regarding etocthenumber in the chapter style definition
   {\makebox[0pt][r]{\textsc{\MakeTextLowercase\etocthenumber}\hspace{\tocsep}}%
    \etocname\kern1em\etocpage\par%
   }%
   {\addvspace{3ex}} % 3ex falls danach Chapter kommt

\etocsetstyle{subsection}
   {\vspace{0ex}}
   {}
   {\makebox[3em][l]{\etocnumber}\etocname\kern1em\etocpage\par}
   {\addvspace{2ex}} % 2ex falls danach Section kommt

\etocsettocstyle{\chapter*{\contentsname}
                \thispagestyle{plain}%
                \leftskip\tocleft\parindent0pt}{}


%----------------------------------------------------------------------------------------
%   OTHER PACKAGES
%----------------------------------------------------------------------------------------
\usepackage{tikz-dependency}  % for dependency trees
\usepackage{tabularx} % for more flexible tables

\usepackage{marginnote} % Enable Notes on the Page Margin
\usepackage{marginfix} % Enables floating for margin figures
\usepackage{ragged2e} % provides better hyphenation, use with camel case: \RaggedRight
\renewcommand*{\raggedleftmarginnote}{\RaggedLeft}
\renewcommand*{\raggedrightmarginnote}{\RaggedRight}

% justified margin notes:
% uncomment the following two lines to for justified layout of margin notes
%\renewcommand*{\raggedleftmarginnote}{\RaggedLeft}
%\renewcommand*{\raggedrightmarginnote}{\RaggedRight}



\renewcommand*{\marginfont}{\setlength{\parskip}{0.5ex}\scriptsize\sffamily} % format margin text



% for sidenotes: change marginpar font
\usepackage{xparse}
\let\oldmarginpar\marginpar
\RenewDocumentCommand{\marginpar}{om}{%
  \IfNoValueTF{#1}
    {\oldmarginpar{\mymparsetup #2}}
    {\oldmarginpar[\mymparsetup #1]{\mymparsetup #2}}}
\newcommand{\mymparsetup}{\scriptsize\sffamily}


% this provides correct alignment for margin text that is inserted
% right at the beginning of a paragraph; however, it messes up the
% alignment in all other cases.
%% therefore, removed for now:
%%\renewcommand{\marginnotevadjust}{0.71\baselineskip}
% The following is the necessary correction for in-paragraph use
\renewcommand{\marginnotevadjust}{0.21\baselineskip}
\renewcommand{\marginnotevadjust}{0.55\baselineskip}

\usepackage{microtype} % enable better typographic setup

\usepackage{multicol} % enable usage of multiple columns

% biblatex setup
% inspired by https://anneurai.net/2017/10/18/thesis-formatting-in-latex/
\usepackage[
  backend=biber,
  style=authoryear-comp,
  maxcitenames=2,
  maxbibnames=4, % when does the citation change to et al?
  uniquelist=false,
  uniquename=false,
  firstinits=true, % author initials in list
  terseinits=true, % no points between initials
  giveninits=true, % always print only initials for given names
  doi=false, % these fields are commonly omitted
  isbn=false, % these fields are commonly omitted
  url=false, 
  dashed=false,
  terseinits=false, % true if no points between initials
  sortcites=true,
  sorting=ynt,
  language=english,
  backref=true, % show on what pages a ref has been cited
]{biblatex} % Use biber backend with alphabetic reference style
\AtEveryBibitem{%
  \clearlist{language} % don't show "en."
  \clearlist{extra} % clears extra fields such as ISBN nrs
}

% shorten the strings used in back references
\DefineBibliographyStrings{english}{%
  backrefpage = {page},
  backrefpages = {pages},
}

%-- no "quotes" around titles of chapters/article titles
\DeclareFieldFormat[article, inbook, incollection, inproceedings, misc, thesis, unpublished]{title}{#1}
%-- no punctuation after volume
\DeclareFieldFormat[article]{volume}{{#1}}
%-- puts number/issue between brackets
\DeclareFieldFormat[article, inbook, incollection, inproceedings, misc, thesis, unpublished]{number}{\mkbibparens{#1}}
%-- and then for articles directly the pages w/o any "pages" or "pp."
\DeclareFieldFormat[article]{pages}{#1}
%-- format 16(4):224--225 for articles
\renewbibmacro*{volume+number+eid}{\printfield{volume}\printfield{number}\printunit{\addcolon}
}
\DeclareFieldFormat{url}{\url{#1}}


\usepackage[autostyle=true]{csquotes} % Required to generate language-dependent quotes in the bibliography

\usepackage[
  obeyFinal,
  textsize=scriptsize,
  backgroundcolor=ubyellow25,linecolor=ubyellow,bordercolor=ubyellow,
]{todonotes}

% change font of todo notes to sans-serif
\makeatletter
\renewcommand{\todo}[2][]{\@bsphack\@todo[#1]{\sffamily #2}\@esphack\ignorespaces}
\makeatother

\usepackage{booktabs} % use formal table layout

\urlstyle{same} % avoids printing URLs in typewriter font

% enable very extensive URL breaking
% https://tex.stackexchange.com/questions/3033/forcing-linebreaks-in-url
\PassOptionsToPackage{hyphens}{url}
\expandafter\def\expandafter\UrlBreaks\expandafter{\UrlBreaks% save the current one
  \do\a\do\b\do\c\do\d\do\e\do\f\do\g\do\h\do\i\do\j%
  \do\k\do\l\do\m\do\n\do\o\do\p\do\q\do\r\do\s\do\t%
  \do\u\do\v\do\w\do\x\do\y\do\z\do\A\do\B\do\C\do\D%
  \do\E\do\F\do\G\do\H\do\I\do\J\do\K\do\L\do\M\do\N%
  \do\O\do\P\do\Q\do\R\do\S\do\T\do\U\do\V\do\W\do\X%
  \do\Y\do\Z\do\*\do\-\do\~\do\'\do\"\do\-}%

% TODO consider using package xurl, which is supposed to handle url breaking

% https://tex.stackexchange.com/a/450695
% allow URLs to be spaced out at / => much better URL breaking in margins
\makeatletter
\g@addto@macro\UrlSpecials
{%
    \do\/{\mbox{\UrlFont/}\hskip 0pt plus 0.1em minus 0.1em}%
}
\Urlmuskip=0mu plus 1mu\relax
\makeatother


% hyperlink layout
\usepackage{hyperref}
 \hypersetup{colorlinks,breaklinks,unicode,
             citecolor=triadicblue,
             linkcolor=triadicblue,
             urlcolor=triadicblue,
             filecolor=triadicblue,}

% cleverref allows you to use \Cref{sec:foo} to get the text "Section 1.2".
% This also works with figures and tables.
\usepackage{cleveref}

% datetime is use to automatically handle the date rendering on the titlepage.
\usepackage{datetime}
% rendering the current date as Month/JJJJ
% see: https://tex.stackexchange.com/questions/212263/month-year-format-in-latex
\newdateformat{monthyeardate}{%
  \monthname[\THEMONTH] \THEYEAR}


\raggedbottom % do NOT force all pages to have the same height (which would be done by increasing the space between paragraphs, which can create noisy layouts)

%----------------------------------------------------------------------------------------
% SETUP BIBLIOGRAPHY
%----------------------------------------------------------------------------------------
\setlength{\bibitemsep}{.3\baselineskip plus .05\baselineskip minus .05\baselineskip}
\newlength{\bibparskip}\setlength{\bibparskip}{0pt}
\let\oldthebibliography\thebibliography
\renewcommand\thebibliography[1]{%
  \oldthebibliography{#1}%
  \setlength{\parskip}{\bibitemsep}%
  \setlength{\itemsep}{\bibparskip}%
}

% allow much more liberal line breaks in URLs
\setcounter{biburllcpenalty}{7000}
\setcounter{biburlucpenalty}{8000}

% adjust space between key and entry, default is 2\labelsep
\setlength{\biblabelsep}{1\labelsep}

% configures indentation of bibentries
\defbibenvironment{bibliography}
  {\list
     {\hspace{0.5\labelalphawidth}\bfseries\printtext[labelalphawidth]{%
        \printfield{prefixnumber}%
        \printfield{labelalpha}%
        \printfield{extraalpha}}}
     {\setlength{\labelsep}{\biblabelsep}%
      \setlength{\leftmargin}{0.5\labelalphawidth}%
      \setlength{\itemsep}{1.5\bibitemsep}%
      \setlength{\parsep}{\bibparsep}}%
      \renewcommand*{\makelabel}[1]{##1\hss}}
  {\endlist}
  {\item}


%----------------------------------------------------------------------------------------
% MARGIN SETTINGS
%----------------------------------------------------------------------------------------

% Using the layout from kaobook
\geometry{
    paper=a4paper,
    head=13.6pt,
    top=27.4mm,
    bottom=27.4mm,
    inner=24.8mm,
    %outer=24.8mm,
    %right=2.183cm,
    textwidth=107mm,
    marginparsep=8.2mm,
    marginparwidth=49.4mm,
    %textheight=49\baselineskip,
    includemp,
    % showframe
}

% Wide figures span text and margin.
% Use the pre-calculated length \widefigurewidth in \includegraphics.
\def\widefigurewidth{\dimexpr(\marginparwidth + \textwidth + \marginparsep)}


%----------------------------------------------------------------------------------------
% SETUP HEADER AND FOOTER
%----------------------------------------------------------------------------------------


\newlength{\overflowingheadlen}
\setlength{\overflowingheadlen}{\textwidth}
\addtolength{\overflowingheadlen}{\marginparsep}
\addtolength{\overflowingheadlen}{\marginparwidth}

% old header/footer, maybe not necessary any more?
\automark[chapter]{chapter}
\ihead{\textup{\headmark}} % Inner header; do not use italics: therefore textup
\ihead{\textup{\textsc{\MakeLowercase\headmark}}}% Inner header - use this line for Small Caps in header
\ohead[]{\pagemark} % Outer header
\cfoot[\pagemark]{} % On chapter opening pages, the page number goes centered into the footer
\automark*[section]{}%

% new header/footer, from kaobook; we could probably remove the original definitions from the cls
\renewpagestyle{thesis}{
  {\hspace{-\marginparwidth}\hspace{-\marginparsep}\makebox[\overflowingheadlen][l]{\textup{\thepage}\quad\rule[-\dp\strutbox]{1pt}{\baselineskip}\quad{}\textup{\textsc{\MakeLowercase \leftmark}}}}%
  {\makebox[\overflowingheadlen][r]{\textup{\textsc{\MakeLowercase \rightmark}}\quad\rule[-\dp\strutbox]{1pt}{\baselineskip}\quad\textup{\thepage}}}%
  {}
}{
  {}%
  {}%
  {}
}
\renewpagestyle{plain.thesis}{
  {}%
  {}%
  {}
}{
  {\hspace{-\marginparwidth}\hspace{-\marginparsep}\makebox[\overflowingheadlen][l]{\textup{\textsc{\thepage}}\quad\rule[-\dp\strutbox]{1pt}{\baselineskip}}}%
  {\makebox[\overflowingheadlen][r]{\rule[-\dp\strutbox]{1pt}{\baselineskip}\quad\textup{\textsc{\thepage}}}}%
  {}
}

%----------------------------------------------------------------------------------------
% LISTINGS SETTINGS
%----------------------------------------------------------------------------------------

\usepackage{textcomp}
\usepackage{listings}
\definecolor{darkgray}{rgb}{.4,.4,.4}

\lstdefinelanguage{JavaScript}{
  keywords={typeof, new, true, false, catch, function, return, null, catch, switch, var, if, in, while, do, else, case, break},
  ndkeywords={class, export, boolean, throw, implements, import, this},
  sensitive=false,
  comment=[l]{//},
  morecomment=[s]{/*}{*/},
  morestring=[b]',
  morestring=[b]"
}

\lstset{
    aboveskip={1\baselineskip},
    abovecaptionskip=-1\baselineskip,
    belowcaptionskip=2ex,
    basicstyle=\footnotesize\ttfamily\linespread{4},
    breaklines=true,
    columns=flexible,
    commentstyle=\color{gray50}\ttfamily\itshape,
    escapechar=@,
    extendedchars=true,
    frame=l,
    framerule=.5pt,
    identifierstyle=\color{black},
    inputencoding=latin1,
    keywordstyle=\color{utred60}\bfseries,
    ndkeywordstyle=\color{utred60}\bfseries,
    numbers=left,
    numbersep=1.25em,
    numberstyle=\scriptsize\ttfamily,
    prebreak = \raisebox{0ex}[0ex][0ex]{\ensuremath{\hookleftarrow}},
    stringstyle=\color{utred60}\ttfamily,
    upquote=true,
    showstringspaces=false,
}

\lstset{literate=%
   *{0}{{{\color{darkgray}0}}}1
    {1}{{{\color{darkgray}1}}}1
    {2}{{{\color{darkgray}2}}}1
    {3}{{{\color{darkgray}3}}}1
    {4}{{{\color{darkgray}4}}}1
    {5}{{{\color{darkgray}5}}}1
    {6}{{{\color{darkgray}6}}}1
    {7}{{{\color{darkgray}7}}}1
    {8}{{{\color{darkgray}8}}}1
    {9}{{{\color{darkgray}9}}}1
}

\lstnewenvironment{latex}
    {\lstset{language=[LaTeX]TeX}}
    {}

%----------------------------------------------------------------------------------------
% MARGINAL CAPTIONS
%----------------------------------------------------------------------------------------

\usepackage{sidenotes}
\usepackage{scrextend} % for ifthispageodd

% objective: instead of having the sidenode number in superscript, we want it like "1:"
\makeatletter
\ExplSyntaxOn
\RenewDocumentCommand \sidenotetext { o o +m }
{
  \IfNoValueOrEmptyTF{#1}
    {
      \@sidenotes@placemarginal{#2}{\thesidenote{}:~#3}
  \refstepcounter{sidenote}
}
    {\@sidenotes@placemarginal{#2}{#1{}:~#3}}
}
\ExplSyntaxOff
\makeatother

% optional objective: automatically justify sidecaptions to match the other marginnotes
% captions of marginfigures etc. shall always be raggedright
% solution from: https://tex.stackexchange.com/questions/358010/subfigures-break-figure-numbering-with-sidecaptions-from-sidenotes-package/358012#358012

\makeatletter
% Instead of "justified" you *can* use "outerraged" in the DeclareCaptionStyle below.
% This may create a inconsistent layout, therefore, we stick to justified by default.
\DeclareCaptionJustification{outerragged}{\ifthispageodd{\RaggedRight}{\RaggedLeft}}

\DeclareCaptionStyle{sidecaption}{format=plain,font={scriptsize,sf},labelfont=bf,margin=0pt,singlelinecheck=true,justification=justified}
\DeclareCaptionStyle{marginfigure}{format=plain,font={scriptsize,sf},labelfont=bf,margin=0pt,singlelinecheck=true}
\DeclareCaptionStyle{margintable}{format=plain,font={scriptsize,sf},labelfont=bf,margin=0pt,singlelinecheck=true}
\DeclareCaptionStyle{widefigure}[justification=centering]{format=plain,font=small,labelfont=bf,justification=RaggedRight,singlelinecheck=true,margin={0px,0px},oneside}
\DeclareCaptionStyle{widetable}[justification=centering]{format=plain,font=small,labelfont=bf,justification=RaggedRight,singlelinecheck=true,margin={0px,0px},oneside}
\makeatother


%----------------------------------------------------------------------------------------
% RESET SIDENOTE COUNTER AT EVERY CHAPTER
%----------------------------------------------------------------------------------------

\let\oldchapter\chapter
\def\chapter{%
  \setcounter{sidenote}{1}%
  \oldchapter
}


%----------------------------------------------------------------------------------------
% SYMBOLS
%----------------------------------------------------------------------------------------

\usepackage{pifont}
\let\oldding\ding% Store old \ding in \oldding
\renewcommand{\ding}[2][1]{\scalebox{#1}{\oldding{#2}}}% Scale \oldding via optional argument
% usage \ding{number} or |ding[factor]{number}


%----------------------------------------------------------------------------------------
% ITEMIZE AND ENUMERATE ENVIRONMENTS
%----------------------------------------------------------------------------------------

\renewcommand{\labelitemi}{\color{utred}{\scalebox{0.8}{\raisebox{0.2ex}{$\blacktriangleright$}}}}
\renewcommand{\labelitemii}{\textbullet}
\usepackage{enumitem}
\setlist[itemize]{parsep=0.8\parskip,left=0pt,topsep=0pt,partopsep=0pt}
\setlist[enumerate]{parsep=0.8\parskip,left=0pt,topsep=0pt,partopsep=0pt}
\setlist[description]{parsep=0.8\parskip,left=0pt,topsep=0pt,partopsep=0pt}


%----------------------------------------------------------------------------------------
% SET PDF METADATA
%----------------------------------------------------------------------------------------

\AtBeginDocument{
\hypersetup{pdftitle=\ttitle} % Set the PDF's title to your title
\hypersetup{pdfauthor=\authorname} % Set the PDF's author to your name
%\hypersetup{pdfkeywords=\keywordnames} % Set the PDF's keywords to your keywords
}
 % Load the settings from Misc/setup.tex
%%%%%%%%%%%%%%%%%%%%%%%%%%%%%%%%%%%%%%%%%%%%%%%%%%%
%
% File: commands.tex
% 
% This file is part of the PSIThesis.cls
% LaTeX documentclass
% 
% The code in this file is made available
% under the following license:
%
% LPPL v1.3c (http://www.latex-project.org/lppl)
%
%%%%%%%%%%%%%%%%%%%%%%%%%%%%%%%%%%%%%%%%%%%%%%%%%%%


% --------------------------------------------------------
% 			CUSTOM COMMANDS FOR BETTER USABILITY
% --------------------------------------------------------


% Custom image command to insert an image
% This command uses 4 required and one optional arguments/parameters with the following meaning:
%
% Optional:
% 1 - Position of the figure (the default position is 't' for top; if no argument is provided, 't' is used
%
% Required:
% 1 - Width of the image
% 2 - Path to the image (inside the figures folder)
% 3 - Caption of the image
% 4 - Label for the image (a universal fig: is prepended)
%
% Required ----------------------------------------------------------------------------
% Optional -------------|		|			  |					|					  |
% 						V		V (1)		  V	(2)				V (3)				  V (4)
% Example Usage: \image[h]{\textwidth}{barplot-before}{This is a fancy barplot.}{barplot-before}
%
% The result will be the same as:
%
% \begin{figure}[h]
% 	\centering
% 	\includegraphics[width=\textwidth]{../figures/barplot-before}
% 	\caption{This is a fancy barplot.}
% 	\label{fig:barplot-before}
% \end{figure}

\newcommand{\image}[5][t]{
	\begin{figure}[#1]
		\centering
		\includegraphics[width=#2]{../figures/#3}
		\caption{#4}
		\label{fig:#5}
	\end{figure}
}

% Custom command to insert two images next to each other with a margin caption
% This command uses 6 required and one optional arguments/parameters with the following meaning:
%
% Optional:
% 1 - Position of the figure (the default position is 't' for top; if no argument is provided, 't' is used
%
% Required:
% 1 - Path to the frist image (inside the figures folder)
% 2 - Path to the second image (inside the figures folder)
% 3 - Caption of the image
% 4 - Label for the image (a universal fig: is prepended)
%
% Required --------------------------------------------------------------------------------------
% Optional -----------------|		  |			    |					|					    |
% 						    V		  V (1)		    V (2)				V (3)				    V (4)
% Example Usage: \twoimages[h]{barplot-before}{barplot-after}{This is a fancy barplot.}{barplot-sidebyside}
%
% The result will be the same as:
%
% \begin{figure}[h]
% 	\centering
% 	\includegraphics[width=0.48\textwidth]{../figures/barplot-before}
% 	\hspace{\fill}
% 	\includegraphics[width=0.48\textwidth]{../figures/barplot-after}
% 	\caption{\label{fig:barplot-sidebyside}This is a fancy barplot.}
% \end{figure}

\newcommand{\twoimages}[5][t]{
	\begin{figure}[#1]
		\centering
		\includegraphics[width=0.48\textwidth]{../figures/#2}
		\hspace{\fill}
		\includegraphics[width=0.48\textwidth]{../figures/#3}
		\sidecaption{\label{fig:#5}#4}[-2\baselineskip]
	\end{figure}
}


% Custom image command for wide figures
% This command uses 4 required and one optional arguments/parameters with the following meaning:
%
% Optional:
% 1 - Position of the figure (the default position is 't' for top; if no argument is provided, 't' is used
%
% Required:
% 1 - Path to the image (inside the figures folder)
% 2 - Caption of the image
% 3 - Label for the image (a universal fig: is prepended)
%
% Required ------------------------------------------------------------------
% Optional -----------------|		  |				 	 |					|
% 						    V		  V (1)		 	  	 V	(2)				V (3)
% Example Usage: \wideimage[h]{barplot-before}{This is a fancy barplot.}{barplot-before}
%
% The result will be the same as:
%
% \begin{figure*}[h]
% 	\includegraphics[width=\widefigurewidth]{../figures/barplot-before}
% 	\caption{This is a fancy barplot.}
% 	\label{fig:barplot-before}
% \end{figure*}
%
% %%%%%%%%%%%%%%%%%%%%%%%
% % 	ATTENTION		%
% %%%%%%%%%%%%%%%%%%%%%%%
% It seems this command messes with the position of other images, therefore it is advised to check the placement of other images.

\newcommand{\wideimage}[5][t]{
	\begin{figure*}[#1]
		\includegraphics[width=\widefigurewidth]{../figures/#2}
		\caption{\label{fig:#4}#3}
	\end{figure*}
}

% This command uses 3 required and one optional arguments/parameters with the following meaning:
%
% Optional:
% 1 - Vertical offset of the figure (the default offset is 1 for one line lower (negative numbers move the image up); if no argument is provided, 1 is used
%
% Required:
% 1 - Path to the image (inside the figures folder)
% 2 - Caption of the image
% 3 - Label for the image (a universal fig: is prepended)
%
% Required ----------------------------------------------------------------------
% Optional -------------------|			|				  |						|
% 							  V			V (1)			  V	(2)					V (3)
% Example Usage: \marginimage[2]{barplot-before}{This is a fancy barplot.}{barplot-before}
%
% The result will be the same as:
%
% \begin{marginfigure}[2\baselineskip]
% 	\includegraphics[width=\marginparwidth]{../figures/barplot-before}
% 	\caption{\label{fig:barplot-before}This is a fancy barplot.}
% \end{marginfigure}
\newcommand{\marginimage}[4][1]{
	\begin{marginfigure}[#1\baselineskip]
		\includegraphics[width=\marginparwidth]{../figures/#2}
		\caption{\label{fig:#4}#3}
	\end{marginfigure}
}
 % Load the custom commands from Misc/commands.tex

% Uncomment this command to make all links black:
%   useful for printing on black-white printers that do a
%   poor job at rasterizing colored text properly
%\hypersetup{colorlinks=false}

\addbibresource{literature.bib} % The filename of the bibliography

%----------------------------------------------------------------------------------------
% THESIS INFORMATION
%----------------------------------------------------------------------------------------
\newcommand{\thesistype}{Bachelor} % type of your thesis (Bachelors, Masters, Doctoral ...)

%%% CHANGE THIS:
% Your thesis title, this is used in the title and abstract, print it elsewhere with \ttitle
\thesistitle{chatconllu: a Two-way Converter between CHILDES and UD Annotations}

% date to be printed on the title, this will automatically update and be in the correct format
% If any changes to this format (Month JJJJ) are necessary the definition can be found in line 337
% of misc/setup.tex
\def\tdate{\monthyeardate\today}

%%% CHANGE THIS:
% Your name, this is used in the title page, print it elsewhere with \authorname
\author{Jingwen Li}

% Your supervisor's name, this is used in the title page, print it elsewhere with \supname
\supervisor{Dr. \c{C}a\u{g}r{\i} \c{C}\"{o}ltekin}

% Your university's name and URL, this is used in the title page, print it elsewhere with \univname
\university{\href{https://uni-tuebingen.de/en/}{Eberhard Karls\\ Universit\"{a}t T\"{u}bingen}}

% Your research group's name and URL, this is used in the title page, print it elsewhere with \groupname
\group{\href{https://uni-tuebingen.de/fakultaeten/philosophische-fakultaet/fachbereiche/neuphilologie/seminar-fuer-sprachwissenschaft/}{Seminar f\"{u}r Sprachwissenschaft}}

% Your department's name and URL, this is used in the title page, print it elsewhere with \deptname
\department{Philosophische Fakult\"{a}t}

% Your faculty's name and URL, this is used in the title page, print it elsewhere with \facname
% TODO: insert *your* degree program in the \faculty command below
% Applied Computer Science
% Computing in the Humanities
% Information Systems
% International Information Systems Management
% International Software Systems Science
% Software Systems Science
% Education in Business and Information Systems
\faculty{Seminar f\"{u}r Sprachwissenschaft\\ \href{https://uni-tuebingen.de/fakultaeten/philosophische-fakultaet/fachbereiche/neuphilologie/seminar-fuer-sprachwissenschaft/}{Philosophische Fakult\"{a}t}}

% Your address, this is not currently used anywhere in the template, print it elsewhere with \addressname
\addresses{address not used}

% Your subject area, this is not currently used anywhere in the template, print it elsewhere with \subjectname
\subject{subject not used}

% Keywords for your thesis, this is not currently used anywhere in the template, print it elsewhere with \keywordnames
\keywords{keywords not used}
%----------------------------------------------------------------------------------------
% END OF THESIS INFORMATION
%----------------------------------------------------------------------------------------


\begin{document}
\hyphenation{CoNLL-U CHAT CLAN MOR GRASP POST}

\selectlanguage{english}

\frenchspacing % do not add additional hspace after end of sentence full stop dot.

\frontmatter % Uses roman page numbering style (i, ii, iii, iv...) for the pre-content pages

\hypersetup{urlcolor=black}

%%%%%%%%%%%%%%%%%%%%%%%%%%%%%%%%%%%%%%%%%%%%%%%%%%%
%
% File: titlepage.tex
% 
% This file is part of the PSIThesis.cls
% LaTeX documentclass
% 
% The code in this file is made available
% under the following license:
%
% LPPL v1.3c (http://www.latex-project.org/lppl)
%
%%%%%%%%%%%%%%%%%%%%%%%%%%%%%%%%%%%%%%%%%%%%%%%%%%%


\pagestyle{plain} % Default to the plain heading style until the thesis style is called for the body content

%----------------------------------------------------------------------------------------
%   TITLE PAGE
%----------------------------------------------------------------------------------------

\begin{titlepage}

\newgeometry{
    inner=4cm, % Inner margin
    outer=4cm, % Outer margin
    marginparwidth=0cm,
    marginparsep=0mm,
    bindingoffset=.5cm, % Binding offset
    top=2.5cm, % Top margin
    bottom=2.5cm, % Bottom margin
    %showframe % Uncomment to show how the type block is set on the page
}

\begin{center}

\vspace*{.06\textheight}

% The logo of University of Bamberg. Do not use the logo without permission.
% Using it on the title page of a thesis is acceptable.
% TODO Add your own logo here. Do yourself and your readers a favor and use
%      a vectorized logo (PDF) instead of a bitmap (PNG, JPEG)
\includegraphics[width=80mm]{misc/UT_WBMW_Rot_4C.pdf}

\vspace*{.06\textheight}

% {\LARGE \textls[130]{\MakeUppercase{\univname}}\par}\vspace{2\baselineskip} % University name

% TODO uncomment the next line when you write a thesis to display the supervisor
% {\large \facname}
% TODO comment out the next line when you write a thesis
% {\large ~} % For use in the guide

% \vspace{1.5cm}

% TODO uncomment the next line when you write a thesis to display the supervisor
\textsc{\Large \thesistype 's~Thesis}\\[1cm] % For use in the thesis
% TODO comment out the next line when you write a thesis
% \textsc{\Large ~}\\[1cm] % For use in the guide

\HRule \\[0.4cm] % Horizontal line
%
{\huge \bfseries \ttitle\par}\vspace{0.3cm} % Thesis title

\HRule \\[0.4cm] % Horizontal line


\textsc{\Large by}\\[0.3cm]

{\large \authorname}

\vspace{2cm}

% TODO uncomment the next two lines when you write a thesis to display the supervisor
\textsc{\Large Supervisor:} \\[0.1cm]
\large \supname\\[2cm]
%%

% \groupname%\\[2cm] % Research group name and department name
{\large \facname}

\vfill

{\large \tdate} % Date

% TODO remove the following lines when you write a thesis
% \vspace{0.25ex}

% {\small \tversion}

% \vspace{0.25ex}

% {\small Links to this document:}\\
% {\small \url{\doiurl}} (initial version)\\
% {\small \url{\githuburl} (most recent version)}
% Remove lines up to this line

\vfill
\end{center}

\restoregeometry

\end{titlepage}


 % Typeset the titlepage

\hypersetup{urlcolor=triadicblue}


%----------------------------------------------------------------------------------------
% QUOTATION
%----------------------------------------------------------------------------------------

% \vspace*{0.2\textheight}

% \noindent\enquote{\itshape Thanks to my solid academic training, today I can write hundreds of words on virtually any topic without possessing a shred of information, which is how I got a good job in journalism.}\bigbreak

% \hfill Dave Barry



%----------------------------------------------------------------------------------------
% ABSTRACT PAGE
%----------------------------------------------------------------------------------------

\begin{abstract}
%\addchaptertocentry{\abstractname}
% uncomment to add the abstract to the table of contents (not recommended)

Both \emph{Child Language Data Exchange System} (CHILDES) and \emph{universal dependencies} (UD) are projects designed to facilitate linguistic research with computerised tools and resources. Each project is paired with a standardised annotation format, which removes barriers to linguistic data-sharing introduced by diverse annotation or transcription habits.

The CHILDES \emph{CHAT} and UD \emph{CoNLL-U} formats, although widely accepted and used in their dominant fields, use different structures and tagsets to represent morphosyntax. With the help of a two-way converter, information can be augmented or cross-validated using existing tools developed for these two annotation formalisms.

To improve the interoperability between CoNLL-U and CHAT, a deterministic mapping for morphosyntactic information from CHAT codes to UD-style tags is created, and \emph{chatconllu} is implemented as a command-line-based two-way converter. Within the limits of a Bachelor thesis, \emph{chatconllu} is developed for and tested on several selected CHILDES databases of English, French and Italian, and the converted files are validated using CHILDES- and UD-maintained tools.

\end{abstract}


%----------------------------------------------------------------------------------------
% ACKNOWLEDGEMENTS
%----------------------------------------------------------------------------------------

%\begin{acknowledgements}
% %\addchaptertocentry{\acknowledgementname}
% Add the acknowledgements to the table of contents (not recommended)
%
%The acknowledgments and the people to thank go here.
%\end{acknowledgements}


%----------------------------------------------------------------------------------------
% TABLE OF CONTENTS
%----------------------------------------------------------------------------------------

\cleardoublepage

% Table of Contents uses a wider layout than the main content
\newgeometry{
		head=13.6pt,
		top=27.4mm,
		bottom=27.4mm,
		inner=24.8mm,
		outer=24.8mm,
		marginparsep=0mm,
		marginparwidth=0mm,
}
{
\hypersetup{linkcolor=black}
\tableofcontents % Prints the ToC entries
}
\restoregeometry

%----------------------------------------------------------------------------------------
% DEDICATION
%----------------------------------------------------------------------------------------

% \dedicatory{For/Dedicated to/To my\ldots}


%----------------------------------------------------------------------------------------
% THESIS CONTENT - CHAPTERS
%----------------------------------------------------------------------------------------
\mainmatter % From here on, numeric (1,2,3...) page numbering
\pagestyle{thesis} % Return the page headers back to the "thesis" style

% Define some commands to keep the formatting separated from the content
\newcommand{\keyword}[1]{\textbf{#1}}
\newcommand{\tabhead}[1]{\textbf{#1}}
\newcommand{\code}[1]{\texttt{#1}}
\newcommand{\file}[1]{\texttt{#1}}
\newcommand{\option}[1]{\texttt{\itshape#1}}

% Figures will automatically be searched for in the Figures subdirectory
\graphicspath{{./figures/}{./examples/}}

%%% CHANGES NEEDED HERE
%
% Include the chapters of the thesis as separate files from the Chapters folder
% Uncomment the lines as you write the chapters
% Mind the \input instead of the \include here, that change is necessary for the appendix formatting
% Due to the \input command you also need to provide the .tex file ending

\chapter{Introduction} % Main chapter title

\label{Chapter1} % For referencing the chapter elsewhere, use \ref{Chapter1}

% intro
Traditionally, linguistic research has primarily been language-specific, with theories stem from studies on Indo-European languages. In recent decades, there has been much progress in crosslinguistic research, one of which is the Universal Dependencies (UD) Project (\cite{nivre2016}, \cite{nivre2020}). As Joakim Nivre explained in his chapter \emph{Towards a Universal Grammar in Natural Language Processing} , instead of a descriptive and explanatory account of morphosyntax, UD's universal grammar is devised for natural language processing (\cite{nivre2015}). Backed by both universal and language-specific tagsets, UD's CoNLL-U annotation is a scheme that allows for crosslinguistic corpus sharing. However, most UD treebanks are still based on written text.

Like CoNLL-U, CHAT (Codes for Human Analysis of Transcripts) (\cite{Macwhinney2000}) is a standardised annotation scheme. It is part of the CHILDES (Child Language Data Exchange System) Project (\cite{Macwhinney2000}), and it offers a standard way to transcribe spontaneous speech. Morphosyntactic information can also be added as dependent tiers to the main utterances, and if preferred, automatically by running MOR (\cite{Macwhinney2000}), POST (\cite{parisse2000}) and GRASP (\cite{Sagae2004}).

It would be interesting to use the advantages of both formalisms to compensate each other. With the help of a converter, CHILDES can provide UD with spontaneous speech treebanks, and UD can offer its tagsets which are widely accepted. There are also many tools developed to process CoNLL-U formatted files, with which information can be augmented to the CHAT files. It is based on this motivation that chatconllu is designed and implemented.

% explain the choice of the 2 annotations first
This thesis aims to explore and experiment with format conversion between these two annotation schemes via designing and implementing a two-way converter with Python. The scope of analysis is limited to a selected few corpora from the CHILDES database. Transcriptions in \emph{CHAT} are first converted to \emph{CoNLL-U} and then converted back to see which information is lost in the process. Both formats are validated using the tools provided on their official websites.

This thesis serves as a complement to the code on GitHub \sidenote{childes-ud: \url{https://github.com/Mentha7/childes-ud}}. \Cref{Chapter2} gives necessary background information of both annotation schemes and outlines related work. \Cref{Chapter3} introduces the converter and its implementation. \Cref{Chapter4} discusses the mappings from CHILDES to UD. \Cref{Chapter5} explains the validation process and discusses the current weaknesses and future improvements of chatconllu. Finally, with \Cref{Chapter6} I conclude this thesis.



% explain the choice of the 2 annotations, to be finished separately and integrated into the section above.
% \section{CHILDES}

% The Child Language Data Exchange System (CHILDES) is part of TalkBank project with focus on child language.

% \section{UD}


\chapter{CHILDES \emph{CHAT} and UD \emph{CoNLL-U}}
\label{Chapter2}

This chapter introduces the two annotation schemes of interest---\emph{CHAT} and \emph{CoNLL-U}---in detail and gives the reader an idea of what information is shared by both annotations and what information is exclusively relevant to only one of them. This will shed light on the implementation of \emph{chatconllu} which will be discussed in \Cref{Chapter3}.

\section{CHILDES CHAT Transcription System}

\subsection{Motivation}
The \emph{Codes for Human Analysis of Transcripts} (CHAT) transcription system is developed for the \emph{Child Language Data Exchange System} (CHILDES) project (\cite{Macwhinney2000}) to aid child language acquisition research.

As MacWhinney pointed out in the CHAT manual, collecting audio data for spontaneous conversational interactions is easy, one just have to turn on the recorder, however, processing the collected recordings and turning them into usable data for linguistic research is a demanding task (\cite{Macwhinney2000}).

For studies that rely on speech corpora, transcribing the conversations is an indispensable preprocessing step. Consistent treatment of raw data is critical to scientific research. Therefore, it is hard to imagine how difficult it was to maintain transcription standards when there was no standard format, especially for child speech, which is more chaotic than average adult speech data. Thus, CHAT is designed to accommodate the need for a standard transcription system of conversational speech. Moreover, CHAT transcripts can be analysed by TalkBank programs \sidenote{\url{https://talkbank.org/software/}} like \emph{Computerized Language Analysis} (commonly referred to as the CLAN Program), which facilitates the transcription, sharing and analysis of human conversational interactions.

\subsection{File Structure}
For a CHAT file to be well-formed, there are some basic structural requirements. The minimal structure of a well-formed CHAT file is described by \emph{minCHAT} (\cite{Macwhinney2000}, pp. 21-22).

% describe only relevant minCHAT structure
First of all, each CHAT file should have at least two parts:
\begin{itemize}
    \item \emph{headers} preceded by the \texttt{@} symbol that stores metadata like language and participants. Some headers are obligatory \sidenote{Obligatory headers include: the file-initial header @Begin, @Languages, @Participants, @ID and the file-final header @End.}, some are optional.\\
    Despite its name, headers can be added in the middle of a CHAT file to represent information like tape locations. It is likely that \emph{header} is a misnomer, and they should be called \emph{meta} as they gernally preserve higher-level information and are not confined to file-inital positions.
    \item the \emph{main line}---utterances led by the asterisk symbol \texttt{*}, which are the actual sentences being said.
\end{itemize}

Then, additional information can be \emph{optionally} added as dependent tiers. Annotations of morphological, syntactical, and phonological information or general remarks of the transcriber's choice all go into separate dependent tiers of the main utterance. Dependent tiers are preceded by the percent symbol \texttt{\%}.

% describe mor and gra tiers
\subsection{Dependent Tiers}
The inclusiveness and extensibility of the CHAT format is demonstrated by the use of dependent tiers, since almost any type of information can be organised into separate dependent tiers and added to the relevant utterances. The CHAT manual defines a standard list of dependent tiers (\cite{Macwhinney2000}, pp. 81-87). The tier names are mostly lowercase three-letter codes, but with the exceptions of tier names preceded by an \texttt{x}, like \texttt{\%xmor}. With the active use of dependent tiers, researchers can adjust the level of specificity in the CHAT files to meet their individual needs. At the same time, since dependent tiers are optional, they do not interfere with other obligatory parts of the file.

A dependent tier can contain one of two types of information:
\begin{itemize}
    \item utterance-level information
    \item word-level information
\end{itemize}

\paragraph{utterance-level information}
Utterance-level information usually appears as descriptive remarks. For example, the action tier \texttt{\%act} that describe actions by the speakers is shown in the example below \sidenote{Excerpt from Adam/040217.cha of the Brown Corpus (\cite{brown1973})}.\\

\lstset{
numbers = none,
frame = single,
}

\begin{lstlisting}[caption={Example of a dependent tier with utterance-level information.}, label={lst:chatsent1}]
*PAU:   0 .
%act:   hits Adam
\end{lstlisting}


\paragraph{word-level information}
Word-level information is associated with words in the utterance. Word-based strings encoded with information are separated by the space character. The below example \sidenote{from the same file in the Brown Corpus as Listing \ref{lst:chatsent1}} shows two dependent tiers with word-level information.

\clearpage

\lstset{
numbers = none,
frame = single,
}

\begin{lstlisting}[caption={Example of dependent tiers with word-level information}, label={lst:chatsent2}]
*MOT:   what happened ?
%mor:   pro:int|what v|happen-PAST ?
%gra:   1|2|SUBJ 2|0|ROOT 3|2|PUNCT
\end{lstlisting}

\paragraph{common tiers}
The two most commonly seen tiers are \texttt{\%mor} and \texttt{\%gra} tiers \sidenote{As shown in Listing \ref{lst:chatsent2}.}, which can be generated automatically by the use of a series of CLAN programs\todo{table: program functions} \sidenote{MOR, POST, POSTMORTEM and GRASP}. The \texttt{\%mor} tier encodes morpgological information including affixes, the stem and part-of-speech category of a word, while the \texttt{\%gra} tier keeps the dependency structure of the utterance and the grammatical relations between head and dependent words.

\newcommand{\conllu}[1]{&\footnotesize\texttt{#1}}
\newcommand{\tab}{&\hspace{0.1em}}
\setlength{\abovedisplayskip}{3pt}
\setlength{\belowdisplayskip}{3pt}
\vspace{-1em}
\section{UD \emph{CoNLL-U} Annotation Scheme}\label{sec:conllu}

The CoNLL-U format (\textcolor{orange}{cite Zeman? here}) is modified from the CoNLL-X format introduced by (\cite{buchholz-marsi-2006-conll}). A file in the CoNLL-U format can be seen as a collection of sentences, where each sentence is followed by an empty line. The empty lines serve as sentence boundaries.

Each CoNLL-U sentence can be divided into two parts:
\begin{itemize}
    \item comment lines \sidenote{Meta-information about the entire document is stored in the comment lines of the first sentence.} start with the hash symbol \texttt{\#} with obligatory fields \texttt{sent\_id} and \texttt{text}. \texttt{sent\_id} corresponds to a unique index in the conllu file and \texttt{text} the text of the sentence. A minimal example is shown below:
    \begin{flalign*}
    \conllu{\# sent\_id = 1}&\\
    \conllu{\# text = why ?}&
    \end{flalign*}
    \item token lines where each line represent one token or punctuation in the sentence.
\end{itemize}
\paragraph{token representation} \sidenote{For more information visit: \url{https://universaldependencies.org/format.html}}
Each token \sidenote{with the exceptions of multi-word tokens} in the text should possess one line, and each line is tab-separated into 10 fields:
\begin{enumerate}
    \item \texttt{ID:} word index starting from 1 for each new sentence \sidenote{With the exception of empty nodes, which can be a decimal number greater than 0.}.
    \item \texttt{FORM:} word form or punctuation as indicated in the sentence comment.
    \item \texttt{LEMMA:} lemma of the word
    \item \texttt{UPOS:} one of the 17 Universal part-of-speech tags \sidenote{The UPOS tags are listed in \Cref{tab:uposset} in \Cref{appendixa}}.
    \item \texttt{XPOS:} language-specific part-of-speech tags
    \item \texttt{FEATS:} a list of morphological features \sidenote{They have to be from UD's universal feature inventory or from a predefined language-specific set accepted by UD. } separated by \texttt{|}.
    \item \texttt{HEAD:} index of the head word of the current token, or \textlf{0} if the current token is the root.
    \item \texttt{DEPREL:} one of the 37 universal syntactic relations in UD \sidenote{The 37 universal syntactic relations can be found on: \url{https://universaldependencies.org/u/dep/index.html}}, or one of its subtypes.
    \item \texttt{DEPS:} head-deprel pair in the form of \texttt{head:deprel}.
    \item \texttt{MISC:} any other information related to the current token.
\end{enumerate}
\paragraph{additional constraints} The fields cannot be empty, if any information \sidenote{with the exception of word index} is not specified for a field, an underscore \texttt{\_} is used in place as a placeholder.

The tokens in the previous minimal example can be represented as follows. The top line in bold is added to indicate the field type:
\begin{align*}
\conllu{\textbf{ID}}\tab\conllu{\textbf{FORM}}\tab\conllu{\textbf{LEMMA}}\tab\conllu{\textbf{UPOS}}\tab\conllu{\textbf{XPOS}}\tab\conllu{\textbf{FEATS}}\tab\conllu{\textbf{HEAD}}\tab\conllu{\textbf{DEPREL}}\tab\conllu{\textbf{DEPS}}\tab\conllu{\textbf{MISC}}\\
\conllu{1}\tab\conllu{why}\tab\conllu{why}\tab\conllu{ADV}\tab\conllu{\_}\tab\conllu{PronType=Int}\tab\conllu{0}\tab\conllu{root}\tab\conllu{\_}\tab\conllu{\_}\\
\conllu{2}\tab\conllu{?}\tab\conllu{?}\tab\conllu{PUNCT}\tab\conllu{\_}\tab\conllu{\_}\tab\conllu{1}\tab\conllu{punct}\tab\conllu{\_}\tab\conllu{\_}
\end{align*}

\paragraph{multi-word tokens}
Multi-word tokens are treated as word groups which span multiple lines. A clear example is given on the UD website and is reproduced on the side to illustrate the point.
\marginnote{
\begin{flalign*}
\conllu{\textbf{ID}}\tab\conllu{\textbf{FORM}}\tab\conllu{\textbf{LEMMA}}\\
\conllu{1-2}\tab\conllu{v\'{a}monos}\tab\conllu{\_}\\
\conllu{1}\tab\conllu{vamos}\tab\conllu{ir}\\
\conllu{2}\tab\conllu{nos}\tab\conllu{nosotros}\\
\conllu{3-4}\tab\conllu{al}\tab\conllu{\_}\\
\conllu{3}\tab\conllu{a}\tab\conllu{a}\\
\conllu{4}\tab\conllu{el}\tab\conllu{el}\\
\conllu{4}\tab\conllu{mar}\tab\conllu{mar}
\end{flalign*}
Example taken from \url{https://universaldependencies.org/format.html}. Fields other than \texttt{ID}, \texttt{FORM} and \texttt{LEMMA} are omitted for simplicity.
}


% subsection  (end)


% \subsubsection{Morphosyntactic Coding with \emph{\%mor} Tier}
% % token links %mor to the main line. one-to-one correspondence between
% % normalised utterance and %mor.
% % structure of individual words
% % structure of word groups  -> conllu: multi-word tokens

% \subsubsection{Syntactic Dependency Analysis with \emph{\%gra} Tier}

\section{Exclusive Information vs. Shared Information}
% token as the basic unit of CHAT and CoNLL-U annotations.
% token also links %mor and %gra to the main line. one-to-one correspondence
% between normalised utterance and %mor.
Understanding how the same type of information is organised and represented in these two formats is crucial to format conversion.

As discussed above, both CHAT and CoNLL-U store meta-information as well as morphosyntactic information of tokens in the utterance. However, upon reading the manuals of both formalisms, it seems that CHAT encompasses a broader range of information types otherwise not present in CoNLL-U formats, like the required headers of a CHAT file. Therefore, although meta data in CHAT is meaningless for CoNLL-U, to reproduce CHAT from a converted file in CoNLL-U format, it must be stored in such a way that it can be put back to the right place.

Since CHAT is more encompassing than CoNLL-U and CoNLL-U files have no exclusive information that require previously non-exist tiers or headers, data contained in CoNLL-U files can be seen as a subset of that of the CHAT files.

The general information types contained in CHAT files are summarised in \Cref{tab:info}. It is clear from the table that file- and sentence-level meta data distributed in diverse places of a CHAT file are all stored as sentence comments in CoNLL-U.
\clearpage

\begin{table}[h!]
\caption {Information organisation in CHAT and CoNLL-U files} \label{tab:info}
\centering
% \begin{tabularx}{\linewidth}{@{}llXl@{}}
\begin{tabularx}{\widefigurewidth}{@{}lcXll@{}}
\toprule
\multicolumn{4}{c}{\textbf{CHAT}}& \multicolumn{1}{c}{\multirow{2}{*}{\textbf{CoNLL-U}}}\\
\cmidrule(lr){1-4}
\textbf{Level}& \textbf{Data type} & \textbf{Information} & \textbf{Storage} &\multicolumn{1}{l}{}\\
\midrule\addlinespace
\multicolumn{1}{l}{\multirow{2}{*}{\textbf{File}}}&\multicolumn{1}{c}{\multirow{2}{*}{Meta}}& obligatory headers    & file begin and end & sentence comments\\
\multicolumn{1}{l}{\multirow{2}{*}{}}&\multicolumn{1}{c}{\multirow{2}{*}{}}& other headers  & middle of the file & sentence comments\\\addlinespace
\cmidrule[0.1pt](lr{1.75em}){1-3}
\addlinespace
\multicolumn{1}{l}{\multirow{5}{*}{\textbf{Sentence}}}&\multicolumn{1}{c}{\multirow{4}{*}{Meta}}& speech codes & main line & sentence comments\\
\multicolumn{1}{l}{\multirow{5}{*}{}}&\multicolumn{1}{c}{\multirow{4}{*}{}}& speaker          & main line & sentence comments\\
\multicolumn{1}{l}{\multirow{5}{*}{}}&\multicolumn{1}{c}{\multirow{4}{*}{}}& dependent tiers  & dependent tiers & sentence comments\\
\multicolumn{1}{l}{\multirow{5}{*}{}}&\multicolumn{1}{c}{\multirow{4}{*}{}}& linkage to media  & main line & sentence comments\\\addlinespace
\multicolumn{1}{l}{\multirow{5}{*}{}}&\multicolumn{1}{c}{\multirow{1}{*}{Text}}& sentence text  & main line & sentence comments\\
\addlinespace
\cmidrule[0.1pt](lr{1.75em}){1-3}
\addlinespace
\multicolumn{1}{l}{\multirow{11}{*}{\textbf{Token}}}  &\multicolumn{1}{c}{\multirow{4}{*}{Morphological}}& word form         & main line & \texttt{FORM}\\
\multicolumn{1}{l}{\multirow{11}{*}{}}&\multicolumn{1}{c}{\multirow{4}{*}{}}& lemma           & \texttt{MOR} stem& \texttt{LEMMA}\\
\multicolumn{1}{l}{\multirow{11}{*}{}}&\multicolumn{1}{c}{\multirow{4}{*}{}}& part-of-speech  & \texttt{MOR} pos& \texttt{UPOS} and \texttt{XPOS}\\
\multicolumn{1}{l}{\multirow{11}{*}{}}&\multicolumn{1}{c}{\multirow{4}{*}{}}& features        & \texttt{MOR} affixes& \texttt{FEATS}\\\addlinespace
\multicolumn{1}{l}{\multirow{11}{*}{}}&\multicolumn{1}{c}{\multirow{3}{*}{Syntactic}}& word index      & \texttt{GRA} index& \texttt{ID}\\
\multicolumn{1}{l}{\multirow{11}{*}{}}&\multicolumn{1}{c}{\multirow{3}{*}{}}& head            & \texttt{GRA} head& \texttt{HEAD}\\
\multicolumn{1}{l}{\multirow{11}{*}{}}&\multicolumn{1}{c}{\multirow{3}{*}{}}& syntactic relation  & \texttt{GRA} GR label& \texttt{DEPREL}\\\addlinespace
\multicolumn{1}{l}{\multirow{11}{*}{}}&\multicolumn{1}{c}{\multirow{4}{*}{Additional}}& clitic type      & \texttt{MOR} & \texttt{MISC}\\
\multicolumn{1}{l}{\multirow{11}{*}{}}&\multicolumn{1}{c}{\multirow{4}{*}{}}& compound components  & \texttt{MOR} & \texttt{MISC}\\
\multicolumn{1}{l}{\multirow{11}{*}{}}&\multicolumn{1}{c}{\multirow{4}{*}{}}& special forms     & main line & \texttt{MISC}\\
\multicolumn{1}{l}{\multirow{11}{*}{}}&\multicolumn{1}{c}{\multirow{4}{*}{}}& ...              & ... & \texttt{MISC}\\\addlinespace
\bottomrule
\end{tabularx}
\end{table}

As for the shared morphosyntactic information, while CHAT uses dependent tiers parallel to the main line, UD takes a token-based approach. Nevertheless, both formalisms use word as the basic unit for morphosyntactic analyses. Thus, although placed in separate tiers, morphological information in the \texttt{\%mor} tier and syntactic relations between words in the \texttt{\%gra} tier can be associated with tokens in the utterance. Similarly, by combining information separated in different fields of the CoNLL-U token line, one can group and reorganise the fields to produce string representations similar to those assigned by MOR and GRASP.

\chapter{Implementation of chatconllu} % Main chapter title

\label{Chapter3} % For referencing the chapter elsewhere, use \ref{Chapter3}

% To make the tool manageable within the frame of a bachelor's thesis, chatconllu is designed based on several selected corpora from the CHILDES database.
\section{Data Structures}
From \Cref{Chapter2} we know that CoNLL-U files are collections of sentences, where each sentence has comment lines and token lines. With \Cref{tab:info}, I also showed that to keep as much information from CHAT files as possible, information exclusively relevant to CHAT had to be stored in sentence comments.

chatconllu takes the idea from the CoNLL-U sentence and proposes to use two object classes---\texttt{Sentence} and \texttt{Token}---to organise the diverse types of information in CHAT files. The benefit of using these two uncomplicated structures is that they correspond well to the CoNLL-U way of information organisation, making writing to CoNLL-U strings straightforward.

\paragraph{\texttt{Sentence}}

Each \texttt{Sentence} object is comprised of a list of \texttt{Token} objects and several additional attributes to store sentence-level information like dependent tiers and file-level metadata like headers preceding the current utterance. The list of \texttt{Token} objects corresponds to the words in the normalised utterance. A normalised utterance is an utterance with the CHAT transcription codes removed. When a \texttt{MOR} tier is present, the tokens should have a one-to-one correspondence with the strings in the \texttt{MOR} tier.
\paragraph{\texttt{Token}}

The attributes of the \texttt{Token} object include the \textlf{10} fields of a normal CoNLL-U token, with two additional attributes:
\begin{itemize}
	\item \texttt{multi}---to store the end index of a multi-word token.
	\item \texttt{type}---to store the clitics type
\end{itemize}

Since \texttt{Token} is organised parallel to the UD-style token, it can be directly converted to a token line \sidenote{Or lines for multi-word tokens}.

\section{Overview of chatconllu}


The behaviour of chatconllu can be described in plain English: For all files or files with name specifications in the given directory, depending on the input file format, convert them to the other format and save the output as different files. I used the \emph{pyconll} \sidenote{pyconll is a low-level API for processing CoNLL-U formatted files.} package developed by (\cite{pyconll}) for parsing the converted CoNLL-U files and accessing the values of token fields. I chose to use it instead of writing my own script because I want to quickly validate the general structure of the converted files. Using the load method from a distributed package allows me to get instant feedback if the file is malformed. chatconllu uses argparse to process inputs from the command line. With optional arguments \sidenote{All commands of chatconllu can be found via the \texttt{chatconllu --help} command after proper installation.}, it's also possible to finetune the desired output by asking chatconllu to generate new dependent tiers or omit existing dependent tiers during format conversion.

% All commands of chatconllu can be found via the \texttt{chatconllu --help} command after proper installation. The optional commands are listed in the Appendix\todo{list optional commands in the Appendix}.

\section{CHAT to CoNLL-U}

The conversion from CHAT to CoNLL-U is outlined as follows:
\begin{enumerate}
	\item parse the CHAT file to obtain utterances grouped with dependent tiers
	\item clean the utterance by removing CHAT transcription codes
	\item extract token information from the utterance-tiers group to construct \texttt{Token} objects
	\item construct \texttt{Sentence} object with the created \texttt{Token} objects and other utterance-level information
	\item write the sentences in CoNLL-U format to an output file
\end{enumerate}

\subsection{Normalising Utterances}\label{sec:normalise}

To clean the utterances, I wrote a helper method \texttt{normalise\_utterances}. This method is a simple parser by itself. The input of this method is an utterance with CHAT transcription codes. The output is a tuple where the first value is a list of tokens in the utterance. After the removal of transcription codes, if the \texttt{\%mor} is given, the length of the list should match the number of strings in the \texttt{\%mor} tier. The second value is simply the line itself, which is to be stored as-is as a sentence comment in the output file, allowing for the recovery of the original utterance in the back-conversion process.


\paragraph{regular expressions}
To remove the CHAT transcription codes, one has to first identify them in the utterance string. Regular expressions are designed for the codes described in the CHAT manual and are tested on the chosen corpora. It is worth noticing that although some regular expressions can be merged into one, they are deliberately left not to, to keep them manageable for future human inspection, addition, editing, and case-wise debugging.

\paragraph{nested angle brackets} The original algorithm I used was simply parsing the string by incrementing the index of the current character and matching the remaining line with defined regular expressions. It was successful in processing flat structures, but not nested structures. Although nested structures are arguably rare in the corpora and the max level of nesting is just one, I still wanted a better solution to the problem.

In CHAT, nested structures occur with \emph{scoped symbols} (\cite{Macwhinney2000}, Part 1, p.72). A scoped symbol is a marked segment of speech. The tokens in the segment are enclosed in angle brackets, while the markers which describe the segment are enclosed in the square brackets that follow. There can be more than one square-bracket-enclosed descriptor following an angle-bracketed speech segment, but no other materials are allowed between the scope enclosed in angle brackets and the symbols enclosed in square brackets.

The challenge of nested angle brackets can be illustrated using two example utterances from the Brown Corpus (\cite{brown1973}):

\texttt{<here <monk(ey)> [/] monkey> [>]}\\  % Sarah/020906 sent 210
\texttt{<<yeah yeah> [/] yeah> [?]} % example from Brown/Sarah/21024, sentence 733

To better understand the examples, I also introduce two operations triggered by the markers:
\begin{itemize}
	\item \emph{deletion}: the previous word or scoped element, as well as this marker, should be removed, triggered by markers like \texttt{[/]}.
	\item \emph{keep}: keep the previous word or the words in the scoped element, remove this marker after that, triggered by markers like \texttt{[?]} and \texttt{[>]}. Note that the \texttt{>} character also appears as a descriptor, making things more complicated.
\end{itemize}

With these two operations defined, it soon becomes clear to the human eye that the two examples should be normalised to \emph{here monkey} and \emph{yeah} after being cleaned. However, the same parses can only be achieved by a computer program using a stack. A stack is used to remember the number of opening angle brackets that have been seen so that the current nesting level is known. With the help of a stack, the first \texttt{>} is correctly associated with the second \texttt{<} in the first sentence and the first {<} in the second example.


% \begin{algorithm}[H]
% \SetAlgoLined\LinesNumberedHidden
% \DontPrintSemicolon
% \KwIn{$object$, $list$, $depth$}
% \KwResult{Push object to list in the level indicated by depth}
% \If{$depth < 0$}{
% 	raise error: mismatch\;
% }
% \While{$depth > 0$}{
% 	$list$ \KwTo last element of $list$\;
% 	$depth$ \KwTo $depth - 1$\;
% }
% $list$.append($object$)
% \end{algorithm}

In chatconllu, nested expressions like the examples above are translated into nested lists with the help of a depth counter and a push operation. The depth counter is like a stack and keeps the current level of nesting, and the push operation pushes the word to the list of the current depth. The algorithm for processing nested structures is listed in Algorithm \ref{alg:nested}. After the utterance string is processed by the \texttt{normalise} method, a nested list corresponding to the structure of the scopes in the utterance is returned. Then, a recursive delete method is used to start removing unwanted elements from the innermost list. Then, the list is flattened and scanned for other operations like omit. The final remaining list is then the list of all tokens to keep.
\clearpage
\SetKwComment{Comment}{\color{green!50!black}// }{}
\renewcommand{\CommentSty}[1]{\textnormal{\ttfamily\color{blue!50!black}#1}\unskip}
\LinesNotNumbered
\DontPrintSemicolon
% \SetKwBlock{kwInit}{Initialization}{end}
\begin{algorithm}[H]\label{alg:nested}
\caption{Dealing with nested angle brackets}
\SetKwFunction{FMain}{Push}
\SetKwProg{Fn}{Function}{:}{}
\Fn{\FMain{$object$, $list$, $depth$}}{
	\If(\tcc*[f]{mismatch}){$depth < 0$}{
		raise error\;
	}
	\While{$depth > 0$}{
		$list$ $\leftarrow$ last element of $list$ \tcc*[f]{go one level deeper}\;
		$depth$ $\leftarrow$ $depth - 1$\;
	}
	$list$.append($object$)
  }
  \;
\SetKwFunction{FMain}{Normalise}
\SetKwProg{Fn}{Function}{:}{}
\Fn{\FMain{$line$}}{
\KwInit{
  $i \leftarrow 0$\;
  $depth \leftarrow 0$\;
  $tokens \leftarrow []$\;
  $chars \leftarrow []$\;
}

	\While{$i < len(line)$}{
		\If{$line[i]$ == <}{
			$push([], tokens, depth)$\;
			$depth \leftarrow depth + 1$\tcc*[f]{stack}\;
			% $i \leftarrow i + 1$\;
		}
		\ElseIf{$line[i]$ == >}{
			$token \leftarrow join(chars)$ \;
			$chars \leftarrow []$ \;
			\If{$token$}{
				$push(token, tokens, depth)$\;
			}
			$depth \leftarrow depth - 1$\tcc*[f]{pop from stack}\;
			% $i \leftarrow i + 1$\;
		}
		...\;
		\Else(\tcc*[f]{reads normal character}){
			$chars \leftarrow line[i]$ \;
		}
		\If(\tcc*[f]{add final token}){$i == len(line)$}{
			$token \leftarrow join(chars)$ \;
			$chars \leftarrow []$ \;
			\If{$token$}{
				$push(token, tokens, depth)$\;
			}
		$i \leftarrow i + 1$\;
		}
	}
	\KwRet $tokens$\;
  }
\end{algorithm}

\paragraph{to normalise or not} Although chatconllu can clean the CHAT annotation codes relatively well, it remains a question whether one really wants to obtain the \emph{clean} version of the utterance. On the one hand, word-based information in \texttt{\%mor} and \texttt{\%gra} tiers is only associated with tokens in the cleaned utterance. Therefore, to use the information provided in these tiers, one has to remove unannotated words (or nonwords). On the other hand, removing things like speech errors makes spontaneous conversational data less speech-like. Or worse, it could make the real sentence structure of the sometimes already ungrammatical utterance even less obvious, at least for dependency parsers. In the study of creating a dependency treebank using Dutch CHILDES corpora, during the process of cleaning the encoded utterances, (\cite{odijk2018anncor}) observed that blindly removing unintelligible words results in a decrease in the performance of the dependency parser they used. For example, although \texttt{xxx} is a nonword, it still functions as a placeholder \sidenote{Think of it as 0 in digits or as a masked token.}. The Alpino parser (\cite{bouma}) they used seems to depend on the positional information, despite the word is unknown. This was again confirmed in a way by (\cite{liu2021}), where they found utterances with omissions are more undesirable in dependency parsing and affected the performance of the dependency parser. Therefore when it comes to preprocessing spoken data like the CHILDES corpora. The degree of normalisation depends on the intended use case.

\subsection{Extracting Information from Dependent Tiers}

Since morphosyntactic information in \texttt{\%mor} and \texttt{\%gra} is kept with individual tokens in CoNLL-U files, they have to be processed in such a way that all relevant information from both tiers (if present) are associated with the respective token in the utterance. It can be more complicated with clitics and contractions, which results in different lengths of the \texttt{\%mor} and \texttt{\%gra} tiers, shown in Listing \ref{lst:chatsent3} \sidenote{from the same file in the Brown Corpus as Listing\ref{lst:chatsent1}}. Clitics, like expressions involving English auxiliary verb \emph{'ve}, the possessive \emph{'s}, and contractions, like the English negation particle \emph{n't} or the French \emph{articles composés}--\emph{du} \sidenote{"du" is the contraction of "de" and "le"}, are treated by the \texttt{MOR} grammar as one word group and assigned a single string with MOR codes, while \texttt{GRASP} gives a separate string representation for each of its parts.\\

\lstset{
numbers = none,
frame = single,
}

\begin{lstlisting}[caption={Example to show the different number of strings in \texttt{\%mor} and \texttt{\&gra} tiers due to contracted form \emph{that's}.}, label={lst:chatsent3}]
*CHI:   dat's [: that's] hard .
%mor:   pro:dem|that~cop|be&3S adv|hard .
%gra:   1|2|SUBJ 2|0|ROOT 3|2|JCT 4|2|PUNCT
\end{lstlisting}

\subsubsection{Morphological information from \%mor}

The syntax of the MOR annotations is complicated by the existence of \emph{word groups}. Word groups are the root cause of the sometimes inconsistent number of annotation strings between \texttt{\%mor} and \texttt{\%gra}. Since the strings in the \texttt{\%mor} tier have a strict one-to-one correspondence with words on the main line, annotations of clitics and contractions like \emph{that's} also have to be single strings.

Here, I introduce the syntax of the MOR annotations for words and word groups, which include clitics and compounds.

\paragraph{word}
The core structure of an individual word is \texttt{pos|stem}, giving only the part-of-speech and stem of the word, but words ususally receive more complex morphological analysis, following the structure below:
$$\texttt{prefix\#pos|stem\&fusionalsuffix-suffix=translation}$$
A word can have any number of prefixes, each followed by a hash \texttt{\#}, fusional suffixes, each preceded by an ampersand \texttt{\&} and suffixes, each preceded by a dash \texttt{-} (\cite{Macwhinney2000}, Part 3, p.7).

\paragraph{word group}
A word group has the word annotation of each word in the word group joined by one of three types of delimiter: the dollar sign \texttt{\$} for preclitics, tilda \texttt{\~{}} for postclitics, and the plus sign \texttt{+} for compounds. Compounds, however, has not only each component word annotation preceded by \texttt{+}, but also receives an additional POS tag which is attached to the beginning of the MOR annotation for the word group. For example, \emph{peanut+butter} in the main line has the MOR annotation \texttt{n|+n|peanut+n|butter}.

% what information needs to be extracted (table)
% \begin{margintable}[1\baselineskip]
% \begin{tabularx}{1\textwidth}{@{}ll@{}}
% \toprule
% \textbf{MOR} & \textbf{Field}\\\midrule
% lemma & \texttt{LEMMA} \\
% part-of-speech & \texttt{UPOS}, \texttt{XPOS}\\
% features & \texttt{FEATS} \\\addlinespace
% clitic type & \texttt{MISC} \\
% compound & \texttt{MISC} \\
% translation & \texttt{MISC} \\\bottomrule
% \end{tabularx}
% \caption{\label{tab:martabmor}Information contained in MOR strings that needs to be extracted, and the corresponding CoNLL-U field to store this information.}
% \footnotesize
% \end{margintable}

Morphological information encoded in the \texttt{\%mor} tier is listed in \Cref{tab:martabmor}. The part-of-speech and features can be extracted directly by parsing the MOR annotation string. They are converted to UD-style values using a deterministic mapping which will be introduced in the next chapter.  During back-conversion from CoNLL-U format, chatconllu keep the original MOR POS tags as language-specific tags in \texttt{XPOS} to reconstruct the MOR annotation. In the future, they should be moved to the MISC field as well. Other information requires more processing and organisation, which I explain below:

\paragraph{lemma} The syntax of the MOR annotations shows that getting the lemma out of a MOR string is not as trivial as it seems. There are several pitfalls to avoid:
\begin{itemize}
	\item The semantics of the symbols may be ambiguous. Take the dash symbol \texttt{-} as an example, for words like \texttt{tic-tac-toe}, the \texttt{-} merely connect parts of the word, but in the example of \texttt{part|get-PRESP} for the word \emph{getting}, \texttt{-} could be an indicator for the present participle suffix \texttt{-ing}.
	\item One should also notice that the stem is not always the lemma, just that for English, these two share the same form in most cases. For instance, for \emph{untied}, we have the MOR segment \texttt{un\#v|tie-PASTP}, which means in plain English that the word has a stem \emph{tie} with prefix \emph{un}, its part-of-speech is verb and this word form is the past participle of the verb. Since in CoNLL-U, the \texttt{LEMMA} field is reserved for lemma and not word stem, the prefix should be put back in place to produce the real lemma.
	\item Compound strings are structured differently and cannot be parsed in the same way. For words like \texttt{n|+n|peanut+n|butter}, instead of parsing this string using a different method, chatconllu simply use the word form as lemma. In this example, the word form and the lemma should both be \texttt{peanutbutter}.
\end{itemize}

\paragraph{misc}
The extensibility of CHAT enables transcribers to note down almost any information they want, but the CoNLL-U format is focused on selective grammatical information of the tokens in the sentences and only values defined either universally or language-specifically by UD are accepted. To protect against the loss of information during format conversion, chatconllu uses the \texttt{MISC} field to store information that is outside of the previous nine fields.
One should also take care not to use reserved symbols like the field-value separator \texttt{|}. For instance, in one wants to save the compound components \texttt{+n|peanut+n|butter} for \texttt{n|+n|peanut+n|butter} as is in , the \texttt{|} character should first be replaced with semantically insignificant symbols like \texttt{@}, and only to be replaced back to \texttt{|} after the \texttt{MISC} field is processed as a dictionary during the reconstruction of dependent tiers.\\
% avoid conflicts replace separator
\vspace{-1.5em}
\subsubsection{syntactic information from \%gra}

% what information needs to be extracted (table)
\begin{margintable}[1\baselineskip]
\begin{tabularx}{1\textwidth}{@{}ll@{}}
\toprule
\textbf{GRA} & \textbf{Field}\\\midrule
word index & \texttt{ID} \\
head & \texttt{HEAD}\\
syntactic relation & \texttt{DEPREL} \\\bottomrule
\end{tabularx}
\caption{\label{tab:martabgra}Information contained in GRA strings that needs to be extracted, and the corresponding CoNLL-U field to store this information.}
\footnotesize
\end{margintable}

The grammatical relation (GR) annotation in \texttt{\%gra} is represented in a simpler form than the annotation in \texttt{\%mor}, as shown in \Cref{tab:martabgra}. And since the annotations in GRA are not affected in form by compounds or multi-word tokens, they can be parsed uniformly. Therefore, I only explain the GRA annotation syntax here for completeness.

A triple, $(\text{dependent}, \text{head}, \text{relation type})$, is used to denote a dependency relation. In the \texttt{\%gra} tier, the annotation of such an relation is of the form \texttt{i|j|g}, where:
\begin{itemize}
	\item \texttt{i} is the index of the current word, which is the dependent
	\item \texttt{j} is the index of the head word
	\item \texttt{g} is the GR label used for this relation
\end{itemize}


\subsection{Sentence-level Information}

Sentence-level info can also be stored in the comments in CoNLL-U. Meta-information given in the headers of CHAT files are saved as file-initial comments before the first sentence because they are file-specific.
However, one has to admit that storing information like this reduces readability. Another thing is that information stored in this way are susceptible to changes accidentally made, which cannot be detected by format validators. One possible solution is to validate the respective format before and after conversion is performed.


\section{CoNLL-U to CHAT}

The back-conversion from CoNLL-U to CHAT is more straightforward than from CHAT to CoNLL-U because CoNLL-U format has a clear structure of organising information. All file- and sentence-level metadata are stored in the comments, and all morphosyntactic information is already categorised and organised into the ten fields of UD token representation discussed in \Cref{sec:conllu}. I describe the processes involved below.

First, with the help of pyconll (\cite{pyconll}), chatconllu loads the CoNLL-U file to obtain pyconll sentences, each of which is composed of comments and tokens. The comments are read as a dictionary with the field names like \texttt{sent\_id} and \texttt{text} being the keys. Then chatconllu opens an output file, extract and write information in the same order as it would appear in a CHAT file, first headers, then the utterances and dependent tiers and finally the last utterances with no tokens and the file-final headers. During the process of writing data to the output file, if \texttt{\%gra} or \texttt{\%mor} is detected, information is reorganised using the syntax of the respective dependent tiers explained earlier. The user can optionally ask chatconllu to generate new empty \texttt{\%gra} or \texttt{\%mor} tiers, with the unspecified values being replaced by an underscore \texttt{\_}. If the converted CoNLL-U files are processed by other tools like UDPipe (\cite{straka-etal-2016-udpipe}) to add or change information, the user can choose to generate new tiers to represent the added information. Both or one of the following new tiers can be added by chatconllu:
\begin{itemize}
	\item \texttt{\%pos}---mirrors the \texttt{\%mor} tier in CHAT but without converting back UPOS tags to the POS tags used by MO
	R. No morphemes are identified by the UD features as well.
	\item \texttt{\%cnl}---corresponds to the \texttt{\%gra} tier. Information also organised as a triplet string of the form \texttt{id|head|deprel}, without converting the values back to GRA codes.
\end{itemize}


\chapter{Deterministic Mapping for Morphosyntax} % Main chapter title
\label{Chapter4}

In the previous chapter, the different ways of how CHAT and CoNLL-U structurally represent morphosyntactic information are explained. This chapter presents a deterministic mapping from CHAT codes to UD-style labels. The mapping is created for English, French and Italian and consists of three major categories: part-of-speech tags, morphological features and dependency relations.

% design differences

% converting MOR codes to UD Universal POS tags
\section{Part-of-Speech Tags}
\label{sec:pos}
% left-aligned table columns with automated line breaks
\newcolumntype{L}{>{\RaggedRight\arraybackslash}X}
\begin{margintable}[1\baselineskip] % move figure down by 1 line
\begin{tabularx}{\textwidth}{@{}Ll@{}}
\toprule
\textbf{MOR} & \textbf{UD}\\ \midrule
adj & ADJ\\
adv & ADV\\
co, fil & INTJ\\
qn, det & DET\\
aux, cop, mod & AUX\\
post, prep & ADP\\
part, inf, neg & PART\\
coord & CCONJ\\
comp & SCONJ\\
num & NUM\\
n & NOUN\\
pro & PRON\\
n:prop & PROPN\\
v & VERB\\
end, beg & PUNCT\\
none, dia & X\\\bottomrule
\end{tabularx}
\caption{\label{tab:martabpos}Examples of MOR POS categories that can be directly converted to UPOS tags. See Table \ref{tab:posmap1} and Table \ref{tab:posmap2} in \Cref{appendixa} for the full mapping.}
\footnotesize
% Examples of MOR POS categories that can be directly converted to UPOS tags. See \ref{tab:posmap1} and \ref{tab:posmap2} in \Cref{appendixa} for the full mapping.
\end{margintable}


According to the CHAT manual, the MOR program, which deals with part-of-speech tagging, has 39 POS tags for English with extensions for other languages, while UD uses a universal POS (UPOS) tagset with 17 tags for all languages in the Universal Dependencies Project (as of version 2.8).

UD's UPOS is developed based on Petrov's \emph{A Universal Part-of-Speech Tagset}, where the crosslinguistic universality of the tagset stems from its operational definition instead of an intrinsic definition of the grammatical categories (\cite{petrov2012}). In CHAT, however, POS tags are much less universal across languages. As described by the CHAT Manual (Section 2.4 Part of Speech Codes), typologically different languages may use quite different tagsets than that prescribed by the MOR program for English. Additionally, researchers can define their own POS tagsets (\cite{Macwhinney2000}), which not only means that in terms of POS tags, CHILDES took an opposite approach from UD, but also makes an exhaustive mapping between POS tags used in CHILDES databases and the UPOS tagset impossible, even just for the three languages chosen. Therefore, the mapping I define here concerns only the standard POS tagset given by the MOR grammar for English, French and Italian.

The adaptation from MOR POS categories to UD's UPOS can be achieved in the following ways:
\begin{itemize}
	\item a direct mapping where the two tags are the same or has a straightforward relationship
	\item collapsing many MOR POS categories into one UPOS tag when applicable
	\item best-effort decision-making based on the description of the POS cetegories by both formalisms
	\item MOR POS categories with no UPOS counterparts are translated as \texttt{X}
\end{itemize}

An example of straightforward mappings between MOR POS codes and UPOS tags is shown in Table \ref{tab:martabpos}. A total of 16 out of the 17 UPOS tags are used in the tagset mapping, with the exception of \texttt{SYM}, which is reserved for symbols.

Collapsing of POS categories can be illustrated by the example of auxiliary verbs. The MOR code differentiates between auxiliary verbs \texttt{aux} or \texttt{v:aux}, copulas \texttt{cop} and modal verbs \texttt{mod}. However, since UD took an operational approach to define grammatical categories, it does not distinguish copulas or modal verbs from other verbs that function as auxiliaries.  In doing so, UPOS can be applied to typologically distant languages.

The following MOR codes require more thoughts during tagset conversion, partially because they do not exist in all languages. For tags like these, UD usually raises them to a more coarse category and specifies the exact type as morphological features.

\paragraph{\texttt{art}}
Articles are ganerally marked as \texttt{DET} in UD Treebanks with feature \texttt{PronType=Art} \sidenote{Definition for \texttt{DET} can be found from \url{https://universaldependencies.org/docs/u/pos/DET.html}}.

\paragraph{\texttt{preart}}
\texttt{preart} are contractions of prepositions and articles, for instance the Italian \emph{alla} is the contraction of \emph{a} and \emph{la}. In UD-style annotations, words like \emph{alla} are treated as multi-word tokens, and a POS tag is given to each part of the token. However, CHAT considers them to be single words with forms of their own. To keep the alignment with \%mor tier, here \texttt{preart} is mapped to \texttt{DET} because the word functions as a determiner.

\paragraph{\texttt{vimp}}
Imperatives are tagged with \texttt{vimp} in \%mor tier. However, in UD, imperatives are tagged with \texttt{VERB}, which is the category for all finite verbforms \sidenote{Definition for \texttt{VERB} can be found from \url{https://universaldependencies.org/docs/u/pos/VERB.html}} . Information of gender, person, number, tense, aspect, mood and voice are all specified as features. Therefore \texttt{vimp}, as an example, is mapped to \texttt{VERB} with feature \texttt{MOOD=Imp}.

\paragraph{\texttt{vpfx}}
Preverb/verb particles are tagged with \texttt{vpfx}. Although they are particles, according to UD's definition \sidenote{Definition for \texttt{ADP} can be found from \url{https://universaldependencies.org/docs/u/pos/ADP.html}}, they belong to \texttt{ADP}.

\paragraph{\texttt{conj}}
Despite the similarity between the MOR \texttt{conj} and UD \texttt{CCONJ}, they represent different things. UD distinguishes between \texttt{CCONJ} \sidenote{Definition for \texttt{CCONJ} can be found from \url{https://universaldependencies.org/docs/u/pos/CCONJ.html}} and \texttt{SCONJ} \sidenote{Definition for \texttt{SCONJ} can be found from \url{https://universaldependencies.org/docs/u/pos/SCONJ.html}}, the former is for coordinating conjunctions like \emph{and} and \emph{but}, which are tagged \texttt{coord} by MOR, and the latter marks subordinating conjunctions with complementizers like \emph{that} or \emph{if}, or adverbial clause introducers like \emph{when}, these are tagged \texttt{conj} by MOR. Therefore, the UPOS tag corresponding to \texttt{conj} is \texttt{SCONJ}.

\paragraph{special form markers}
It is worth noticing that some MOR codes do not correspond to any grammatical categories. For example, in CHAT, the so-called \emph{special form markers} have unique POS tags independent of the part-of-speech category to which the word belongs. Nevertheless, some of these tags can still be mapped to UPOS, like \texttt{sing} for singing, \texttt{bab} for babbling and \texttt{wplay} for wordplay can all be mapped to \texttt{INTJ}, which stands for \emph{interjection}. The definition of interjection, however, is vague. UD defines interjection descriptively \sidenote{"An interjection is a word that is used most often as an exclamation or part of an exclamation. It typically expresses an emotional reaction, is not syntactically related to other accompanying expressions, and may include a combination of sounds not otherwise found in the language." Excerpt from \url{https://universaldependencies.org/docs/u/pos/SCONJ.html}}. The core criteria for a word \sidenote{or non-word in this case} to be an interjection include:
\begin{itemize}
	\item that it is often used exclamatorily
	\item that it typically expresses emotions
	\item that it plays no syntactic role in the utterance
	\item that it might not be a recognised word in the language
\end{itemize}
Since singing, babbling and wordplay satisfy these criteria, along with the MOR code explicitly for interjections \texttt{co}, they are granted the \texttt{INTJ} tag. Other special form markers, on the other hand, has no corresponding UPOS tags whatsoever, for example \texttt{dia} for dialect forms. Unfortunately they have to be assigned to \texttt{X}, for words that does not belong to the other 16 categories.

\begin{margintable}[1\baselineskip] % move figure down by 1 line
\begin{tabularx}{0.85\textwidth}{@{}cl@{}}
\toprule
Punctuations & MOR Code\\\midrule
„ & end\\
‡ & beg\\
, & cm\\
“ & bq\\
” & eq\\
‘ & bq2\\
’ & eq2\\\bottomrule
\end{tabularx}
\caption{\label{tab:martabpunct}MOR punctuation marks and their corresponding MOR codes.}
\footnotesize
\end{margintable}

\paragraph{punctuation marks}
Finally, MOR also defines a set of punctuation marks with string representations, shown in Table \ref{tab:martabpunct}. Although treated as lexical items in the \%mor tier, they are all given the \texttt{PUNCT} tag.

% feature mappings
\section{Morphological Features}
\label{sec:feats}

Morphological features like \emph{tense} and \emph{mood} are encoded in the string representation of the word. MOR uses a morpheme-based approach to analyse individual words. As discussed in the previous chapter, a word is separated into prefix, part-of-speech, stem, affixes or morphemes indicating specific morphological features. Unlike MOR's dissecting analysis of the word, UD treats words as the basic unit, and features are considered properties of the word. It does not presuppose the existence of morphemes.

\begin{margintable}[1\baselineskip]
\begin{tabularx}{1\textwidth}{@{}ll@{}}
\toprule
\textbf{CHILDES Database} & \textbf{Language}\\ \midrule
	Brown & eng\\
	English-MiamiBiling & eng\\
	Belfast & eng\\
	Wells & eng\\
	MTLN & fra\\
	Leveille & fra\\
	vioncolas & fra\\
	Geneva & fra\\
	Tonelli & ita\\\bottomrule
\end{tabularx}
\caption{\label{tab:martabdb}CHILDES databases chosen for this project.}
\footnotesize
\end{margintable}

These two fundamentally different views of words mean that morphological information is stored in different ways. CoNLL-U format keeps all relevant features of a word in the token's \texttt{feats} field. And although each language in UD has its subset of morphological features, the feature types and values are by convention the same. Compared with UD, CHILDES's way of handling features is less direct. Not all features are explicit as UD's feature-value pairs. Instead, they have to be inferred from the morpheme codes or POS tags.

Listing of all possible MOR codes for grammatical morphemes is, again, unlikely. Therefore, I chose the easier path of collecting these codes from the CHILDES databases that I decided to work with. Table \ref{tab:martabdb} lists the corpora of my choice. Of all morpheme codes found in these corpora, not all of them correspond to UD features. Since morphological features are not mandatory, I only created mappings for codes with a corresponding UD-style feature-value pair. An example mapping is given in Table \ref{tab:martabfeats} on the right-hand side. The features that can be obtained from such codes are limited to 9 types: \texttt{Number}, \texttt{Person}, \texttt{Gender}, \texttt{Mood}, \texttt{Degree}, \texttt{Tense}, \texttt{Verbform}, \texttt{Voice}, and \texttt{Poss}. All possible UD features and values allowed for Engish, French, and Italian can be found in Table \ref{tab:featsmap} in \Cref{appendixa}. Notice that since the codes are sometimes capitalised and sometimes not, they are all converted to lowercase for easy comparison. The original morpheme codes are kept in the \texttt{MISC} field, which keeps all token-level information that does not fit into the previous nine fields.

\begin{margintable}[1\baselineskip]
\begin{tabularx}{1\textwidth}{@{}ll@{}}
\toprule
\textbf{MOR} & \textbf{feats}\\ \midrule
	3s & Number=Sing|Person=3\\
	cp & Degree=Cmp\\
	f & Gender=Fem\\
	inf & VerbForm=Inf\\
	imp & Mood=Imp\\
	impf & Tense=Imp\\
	pass & Voice=Pass\\
	pastp & Tense=Past|VerbForm=Part\\
	poss & Poss=Yes\\\bottomrule
\end{tabularx}
\caption{\label{tab:martabfeats}Example MOR grammatical morpheme codes and their corresponding UD featue-value pairs.}
\footnotesize
\end{margintable}


% converting GRA codes to UD deprels
\section{Dependency Relations}
\label{sec:deprel}

In CHAT annotations, syntactic relations are represented by Grammatical Relations (GRs), first devised by Kenji Sagae (\cite{sagae-etal-2004-adding}) for the English language. After later revisions in 2007, CHAT now has 37 dependency relations (\cite{sagae2007}). They correspond to \emph{deprels} (dependency relations) in CoNLL-U annotations. However, CHILDES and UD use different labels for these dependency relations and quite different parsing schemes. Sometimes, a conversion from GR to deprel also introduces a change in the structure of the dependency tree.

Defining mappings between a GR and a UD deprel thus ranges from direct label translation to linguistic knowledge-based conditional decisions. Direct mappings are listed in Table \ref{tab:martabgr}. Explanations of the adaptations made will be elaborated below and with detailed examples for a selected few hard-to-decide cases.

\begin{margintable}[1\baselineskip]
\begin{tabularx}{1\textwidth}{@{}Ll@{}}
\toprule
\textbf{GR} & \textbf{deprel}\\ \midrule
SUBJ & nsubj\\
CSUBJ & csubj\\
OBJ & obj\\
OBJ2 & iobj\\
APP & appos\\
AUX & aux\\
CONJ & conj\\
COORD & cc\\
DET & det\\
PUNCT & punct\\
ROOT & root\\\bottomrule
\end{tabularx}
\caption{\label{tab:martabgr}Example CHILDES grammatical relations and their corresponding UD dependency relations.}
\footnotesize
\end{margintable}

\paragraph{\texttt{INCROOT}}
Since CHAT transcripts authentically reflect conversational speech, the utterances may not be fully grammatical. This ungrammaticality problem becomes even more apparent when dealing with child speech. For incomplete utterances or sentences missing the main verb, a substitute root \texttt{INCROOT} is given by GRASP. However, most of the time, it does not hurt to give it the regular \texttt{ROOT} label.
\clearpage

\begin{minipage}{\widefigurewidth}
\begin{dependency}[edge slant=3pt]
	\begin{deptext}[column sep=0.7cm]
	well \& ‡ \& that \& 's \& life \& .\\
	% \textlf{1}|\textlf{0}|BEG \& \textlf{2}|\textlf{1}|BEGP \& \textlf{3}|\textlf{4}|SUBJ \& \textlf{4}|\textlf{0}|ROOT \& \textlf{5}|\textlf{4}|OBJ \& \textlf{6}|\textlf{4}|PUNCT\\
	\end{deptext}
	\depedge{1}{2}{BEGP}
	\depedge{4}{3}{SUBJ}
	\depedge{4}{5}{OBJ}
	\depedge{4}{6}{PUNCT}
	\deproot{1}{BEG}
	\deproot{4}{ROOT}
\end{dependency}
\hfill
\begin{dependency}[edge slant=3pt]
	\begin{deptext}[column sep=0.7cm]
	well \& , \& that \& 's \& life \& .\\
	\end{deptext}
	\depedge{5}{1}{discourse}
	\depedge{5}{2}{punct}
	\depedge{5}{3}{nsubj}
	\depedge{5}{4}{cop}
	\depedge{5}{6}{punct}
	\deproot[edge unit distance=4ex]{5}{root}
\end{dependency}
\end{minipage}
\captionof{figure}{CHILDES vs UD dependency graphs: interjections}\label{fig:intj}

\begin{minipage}{\widefigurewidth}
\begin{dependency}[edge slant=3pt]
	\begin{deptext}[column sep=0.7cm]
	son \& ‡ \& that \& 's \& life \& .\\
	% \textlf{1}|\textlf{0}|BEG \& \textlf{2}|\textlf{1}|BEGP \& \textlf{3}|\textlf{4}|SUBJ \& \textlf{4}|\textlf{0}|ROOT \& \textlf{5}|\textlf{4}|OBJ \& \textlf{6}|\textlf{4}|PUNCT\\
	\end{deptext}
	\depedge{1}{2}{BEGP}
	\depedge{4}{3}{SUBJ}
	\depedge{4}{5}{OBJ}
	\depedge{4}{6}{PUNCT}
	\deproot{1}{BEG}
	\deproot{4}{ROOT}
\end{dependency}
\hfill
\begin{dependency}[edge slant=3pt]
	\begin{deptext}[column sep=0.7cm]
	son \& , \& that \& 's \& life \& .\\
	\end{deptext}
	\depedge{5}{1}{vocative}
	\depedge{5}{2}{punct}
	\depedge{5}{3}{nsubj}
	\depedge{5}{4}{cop}
	\depedge{5}{6}{punct}
	\deproot[edge unit distance=4ex]{5}{root}
\end{dependency}
\end{minipage}
\captionof{figure}{CHILDES vs UD dependency graphs: vocatives}\label{fig:voc}

\subsection{Difficult Cases}\label{sec:grdifficult}
\subsubsection{The multiple-root problem}
In UD, every sentence has exactly one root, resulting in a single-root dependency tree. The dependency analysis generated by GRASP is, however, not always like that.\\

Sentences with non-final punctuations, like after an exclamation, a vocative expression or, close to the end of the utterance, before the tag question, after processed by GRASP, receives multiple roots. Although usually only one \texttt{ROOT} or \texttt{INCROOT} label is assigned, there are multiple words whose head is marked \textlf{0}. This, in effect, corrupts the single-root tree structure.\\

This problem occurs because exclamations and vocatives are not considered part of the clause. Tag markers are, however, for the chosen corpora, not so much of a problem and usually have their head marked to the functional root (\texttt{ROOT} or \texttt{INCROOT}), like in \Cref{fig:taggr}.\\

Consider the following sentences:

\pex~ Sentences with non-final punctuations\label{roots}
\a {[\sl Well\/}]$_{\text{\sc intj}}$, that's life.        \hfill {\sl interjection}\label{intj}
\a {[\sl  Son\/}]$_{\text{\sc voc}}$, that's life.         \hfill {\sl vocative}\label{voc}
\a That's life , {[\sl isn't it}]$_{\text{\sc tag}}$?    \hfill {\sl tag question}\label{tagq}
\xe


After running GRASP, sentences like (\ref{intj}) and (\ref{voc}) will be analysed like the left-hand side dependency graph of \Cref{fig:intj} and \Cref{fig:voc}, while UD-style dependeny analyses are shown on the right-hand side.\\


It is clear from the dependency graphs that the CHILDES trees are broken. Therefore, to produce a dependency tree from the pieces, one must define a change-head operation, which changes the token's head index to that of the root of the sentence.\\


After the change-head operation is applied, the current token with GR \texttt{BEG} should be connected to \texttt{ROOT}. Now the two examples in \Cref{fig:intj} and \Cref{fig:voc} share the same dependency structures, but the dependency relation labels should be different for vocatives and interjections. chatconllu needs to know the POS tags of the token with \texttt{BEG} as GR to be able to make the final decision. The decision process is outlined below:\\

\begin{itemize}
	\item If the token has a POS tag and it is not \texttt{INTJ}, \texttt{PROPN}, or \texttt{NOUN}, the corresponding \texttt{deprel} \sidenote{Definition for \texttt{parataxis} can be found from \url{https://universaldependencies.org/u/dep/parataxis.html}} label should be \texttt{parataxis} for a loose side-by-side placement between the elements.
	\item If the POS tag is \texttt{INTJ}, the corresponding \texttt{deprel} label should be \texttt{discourse} \sidenote{Definition for \texttt{discourse} can be found from \url{https://universaldependencies.org/u/dep/discourse.html}}.
	\item If the POS tag is \texttt{PROPN} or \texttt{NOUN}, the corresponding \texttt{deprel} label should be \texttt{vocative} \sidenote{Definition for \texttt{vocative} can be found from \url{https://universaldependencies.org/u/dep/vocative.html}}.
\end{itemize}

% For instance, the sentence "well, son, that's just life, isn't it ?" will be analysed by GRASP as this:
\subsubsection{Inconsistent Dependency Parses}

% \begin{minipage}{\widefigurewidth}
\begin{minipage}{\widefigurewidth}
\begin{dependency}[edge slant=3pt]
	\begin{deptext}[column sep=1cm]
	that \& 's \& life \& „ \& is \& n't \& it \& ?\\
	% \textlf{1}|\textlf{2}|SUBJ \& \textlf{2}|\textlf{0}|ROOT \& \textlf{3}|\textlf{2}|OBJ \& \textlf{4}|\textlf{5}|ENDP \& \textlf{5}|\textlf{2}|END \& \textlf{6}|\textlf{2}|PUNCT\\
	\end{deptext}
	\depedge{2}{1}{SUBJ}
	\depedge{2}{3}{PRED}
	\depedge{5}{4}{ENDP}
	\depedge[edge unit distance=1.8ex]{2}{5}{END}
	\depedge{5}{6}{NEG}
	\depedge[edge unit distance=2.5ex]{5}{7}{PRED}
	\depedge[edge unit distance=1.5ex]{2}{8}{PUNCT}
	\deproot[edge unit distance=4ex]{2}{ROOT}
\end{dependency}
% \captionof{figure}{CHILDES}
\end{minipage}
\captionof{figure}{CHILDES: tag questions}
\label{fig:taggr}
\vspace*{2cm}
% \begin{minipage}{\widefigurewidth}\label{fig:tagdeprel}
\begin{minipage}{\linewidth}
% \begin{dependency}[arc edge, arc angle=70, text only label, label style={above}]
\begin{dependency}[edge slant=3pt]
	\begin{deptext}[column sep=1cm]
	that \& 's \& life \& , \& is \& n't \& it \& ?\\
	\end{deptext}
	\depedge{3}{1}{nsubj}
	\depedge{3}{2}{cop}
	\depedge{3}{4}{punct}
	\depedge{3}{5}{parataxis}
	\depedge{5}{6}{advmod}
	\depedge{5}{7}{nsubj}
	\depedge[edge unit distance=1.8ex]{3}{8}{punct}
	\deproot[edge unit distance=3ex]{3}{root}
\end{dependency}
% \captionof{figure}{UD}
\end{minipage}
\captionof{figure}{UD: tag questions}
\label{fig:tagdeprel}

\clearpage
\begin{minipage}[b]{0.5\linewidth}
\begin{dependency}
	\begin{deptext}[column sep=0.5cm]
	She \& is \& a \& student \& .\\
	% PRON \& AUX \& DET \& NOUN \& PUNCT\\
	\end{deptext}
	\depedge{4}{1}{nsubj}
	\depedge{4}{2}{aux}
	\depedge{4}{3}{det}
	\depedge{4}{5}{punct}
	\deproot{4}{root}
\end{dependency}
\end{minipage}
\hfill
\begin{minipage}[b]{0.5\linewidth}
\begin{dependency}
	\begin{deptext}[column sep=0.5cm]
	She \& is \& a \& student \& .\\
	% PRON \& AUX \& DET \& NOUN \& PUNCT\\
	\end{deptext}
	\depedge{2}{1}{SUBJ}
	\depedge{2}{4}{PRED}
	\depedge{4}{3}{DET}
	\depedge{2}{5}{PUNCT}
	\deproot{2}{ROOT}
\end{dependency}
\end{minipage}

\begin{minipage}[b]{0.5\linewidth}
\begin{dependency}
	\begin{deptext}[column sep=0.5cm]
	She \& became \& happier \& .\\
	% PRON \& VERB \& ADJ \& PUNCT\\
	\end{deptext}
	\depedge{2}{1}{nsubj}
	\depedge{2}{3}{xcomp}
	\depedge{2}{5}{punct}
	\deproot{2}{root}
\end{dependency}
\end{minipage}
\begin{minipage}[b]{0.5\linewidth}
\begin{dependency}
	\begin{deptext}[column sep=0.5cm]
	She \& became \& happier \& .\\
	% PRON \& VERB \& ADJ \& PUNCT\\
	\end{deptext}
	\depedge{2}{1}{SUBJ}
	\depedge{2}{3}{PRED}
	\depedge{2}{5}{PUNCT}
	\deproot{2}{ROOT}
\end{dependency}
\end{minipage}

\chapter{Validation and Comparison} % Main chapter title

\label{Chapter5} % For referencing the chapter elsewhere, use \ref{Chapter5}

\section{Validating CHAT using CLAN Tools}

After using chatconllu to convert the CoNLL-U files back to CHAT format, the user may want to validate that the resulting CHAT files are well-formed and accepted by the CLAN program. CLAN has a CHECK program \sidenote{CLAN can be downloaded via \url{https://dali.talkbank.org/clan/}} that can validate the CHAT transcription format. Alternatively, if the user wants to stick to the command line, they can download Chatter \sidenote{Chatter - CHAT format validator: \url{https://talkbank.org/software/chatter.html}}, a Java program from TalkBank, and follow the instructions on the website.\\

% include a use case.

For the chosen corpora, CHAT files produced by chatconllu pass the CHECK program in CLAN but not Chatter for the reason of missing tabs. Long lines in the original cha files generated from the CLAN program are broken into separate lines after a certain amount of characters is reached. Continuations of the line are preceded by an initial tab. However, the current version of chatconllu does not handle the maximum length of a line, and because of that, it fails the Chatter validation process.\\

\section{Validating CoNLL-U using UD Tools}

UD also provides tools for format validation. One of the UD maintained tools is the script \emph{validate.py} \sidenote{link to file on GitHub: \url{https://github.com/UniversalDependencies/tools/blob/master/validate.py}}. Its validation are in five levels. To run the script, one has to input a language code because in high-level tests, this script also check whether the values in the fields are allowed for that language. Currently, chatconllu only pass the first two levels at best, which does not check if the values are accepted by UD or specifications of the given language.

\section{Comparison with Original Files}
The differences between the original and back-converted cha files can be visually inspected using \emph{icdiff}. With this command, a side-by-side display of the two cha files selected is displayed in the terminal, with the differences coloured. As mentioned above, most visible differences are the results of the missing leading tabs. Unfortunately, I haven't figured out a way to solve this problem programmatically.

\section{Related Work}

\paragraph{Dependency Parsing for Low-resource Spontaneous Speech} % (fold)
\label{par:zoey}
(\cite{liu2021}) used a low-resource setting, focusing on one child from the Brown Corpus (\cite{brown1973}) and tested the performance of dependency parsers on parent-child conversations. The out-domain parser was trained on written text and performed well on adult speech but not so well on child speech, while the parser trained on a limited amount of in-domain spoken data improved the parsing result of child speech to be comparable with adult speech. As part of their study, they created a semi-automatic adaptation of annotation formats from CHILDES to UD.

\paragraph{The AnnCor CHILDES Treebank}
(\cite{odijk2018anncor}) preprocessed the Dutch CHILDES corpora, augmented them with syntactic annotations using the Alpino parser, and performed partial manual verification. While creating this treebank, they noticed the many disadvantages of annotation discrepancies and argued that annotation guidelines should be explicit and should be strictly followed to produce consistency. They suggest combining human annotations with automatic checks by computer programs. They also pointed out that although the grammar used in child speech is limited since the child is still acquiring the language, assigning adult syntactic structures to their utterances is still valid.

\todo{add comments}

\section{Future Improvements}
\todo{add areas of improvements}

\chapter{Conclusion} % Main chapter title

\label{Chapter6} % For referencing the chapter elsewhere, use \ref{Chapter6}

This thesis project implements a two-way converter, chatconllu, as a command-line tool written in Python for CHILDES CHAT and UD CoNLL-U format. Conversion from CHAT to CoNLL-U cleans the utterances with CHAT transcription codes, extracts information from the dependent tiers, associates them with their corresponding words in the cleaned sentences, and writes the sentences in CoNLL-U format. Back-conversion from CoNLL-U files outputs the original CHAT files, but the tabs for line continuations are missed during this process. Additionally, chatconllu can be used as an intermediate step to augment information to either annotation format. For instance, one can start with converting CHAT files to CoNLL-U format and pass them to any type of dependency parsers that work with this format. After that, the files can either be exported as a preliminary spoken treebank for human inspection, or converted back to CHAT format with added dependent tiers when asked to do so.\\

Moreover, as an attempt, a deterministic mapping to transform CHILDES morphosyntactic codes to UD values is created. Currently, this mapping oversimplifies the problem of tagset conversion by ignoring many non-deterministic cases and assigning uniform values, even though they can be decided using rule-based procedures.\\

In the future, chatconllu can be extended to support visualisations of dependency trees, which will ease the process of designing rules for tagset conversions by helping users experienced in dependency annotations to identify mal-formed and atypical dependency trees. Furthermore, in the hope of aiding independent researchers to tailor the program according to their needs, chatconllu can also be improved by giving the user more control over the output content.

% \chapter{Ideas} % Main chapter title

\label{Chapter7} % For referencing the chapter elsewhere, use \ref{Chapter7}

% motivation?
Unlike most written text or lab speech data, spontaneous speech is <chaotic> in form. Without drafting or editing, a sentence settle in form as each word is spoken. Errors and markers of disfluency like pauses and repetitions are ubiquitous in spontaneous speech data. Before a consensus is reached in how people transcribe and annotate speech, transcribers had to device their own markings/symbols to faithfully transcribe the utterances/conversational interactions. The disadvantages are clear, inconsistent use of symbols is an obstacle for the sharing of linguistic data, which is indispensable in Natural Language Processing.\\


\section{Related Work}

% references
CHILDES database \cite{Macwhinney2000}\\
the Brown Corpus <needs bibtex (Brown, 1973)>\\
UPOS <Petrov et al., 2012>
UDPipe <Straka et al., >
UD Project <Nivre et el., 2016>

\section{Deterministic Mapping for Morphosyntax} % (fold)
\label{sec:mappings}
In the previous section, the different ways of how CHAT and CoNLL-U structurally represent morphosyntactic information are explained. This chapter presents a deterministic mapping from CHAT codes to UD-style labels. The mapping is created for English, French and Italian and consists of three major categories: part-of-speech tags, morphological features and dependency relations.

% design differences

% converting MOR codes to UD Universal POS tags
\subsection{Part-of-Speech Tags}

According to the CHAT manual, the MOR program, which deals with part-of-speech tagging, has 39 POS tags for English with extensions for other languages, while UD uses a universal POS (UPOS) tagset with 17 tags for all languages in the Universal Dependencies Project (as of version 2.8). UD's UPOS is developed based on Petrov's "A Universal Part-of-Speech Tagset", where the crosslinguistic universality of the tagset stems from its operational definition instead of an intrinsic one (Petrov, 2012). In CHAT, however, POS tags are much less universal across languages. As described by the CHAT Manual (Section 2.4 Part of Speech Codes), typologically different languages may use a very different tagset than that prescribed by the MOR program for English. Additionally, researchers can define their own POS tagsets (MacWhinney), which makes an exhaustive mapping between POS tags used in CHILDES databases and the UPOS tagset impossible, even just for the three languages chosen. Therefore, the mapping I define here concerns only the standard POS tagset given by the MOR lexicon files for English, French and Italian.

The adaptation from MOR POS categories to UD's UPOS can be achieved in the following ways:
\begin{itemize}
	\item a direct mapping where the two tags are the same or has a straightforward relationship
	\item collapsing many MOR POS categories into one UPOS tag
	\item best-effort decision-making based on the description of the POS cetegories by both formalisms
	\item MOR POS categories with no UPOS counterparts are translated as \texttt{X}
\end{itemize}

\begin{table}[h!]
\begin{tabular}{@{}lllll@{}}
\toprule
\textbf{MOR POS} & \textbf{UD UPOS} & \textbf{Description}\\ \midrule
adj, adj:pred & ADJ & adjectives\\
adv, adv:tem, neg & ADV & \\
co, fil, wplay, bab, sing & INTJ & \\
qn, det, quant, art, prepart & DET & \\
aux, v:aux, cop, mod & AUX & \\
post, prep, vpfx & ADP & \\
part, inf & PART & \\
coord & CCONJ & \\
conj, comp & SCONJ & \\
num & NUM & \\
n, on, onoma & NOUN & \\
pro & PRON & \\
n:prop & PROPN & \\
v, vimp & VERB & \\
end, beg, cm, bq, eq, bq2, eq2 & PUNCT & punctuation marks\\
none, dia, test, meta, phon, fam, uni, L2, neo, chi & X & \\\bottomrule
\end{tabular}
\end{table}

\paragraph{\texttt{art}}
Articles are marked as \texttt{DET} in UD Treebanks with feature \texttt{PronType=Art}.\\
\paragraph{\texttt{preart}}
\texttt{preart} are contractions of prepositions and articles, for instance the Italian \emph{alla} is the contraction of \emph{a} and \emph{la}. In UD-style annotations, words like \emph{alla} are treated as multi-word tokens and a POS tag is given to each part of the token. However, CHAT consider them to be single words with forms of their own. To keep the alignment with \%mor tier, here \texttt{preart} is mapped to \texttt{DET} because the word functions as a determiner.\\

\paragraph{\texttt{vimp}}
Imperatives are tagged with \texttt{vimp} in \%mor tier. However, in UD, imperatives are tagged with \texttt{VERB}, which is the category for all finite verbforms. Information of gender, person, number, tense, aspect, mood and voice are all specified as features. Therefore \texttt{vimp}, as an example, is mapped to \texttt{VERB} with \texttt{MOOD=Imp}.\\

\paragraph{\texttt{vpfx}}
Preverb/verb particles are tagged with \texttt{vpfx}. Although they are particles, according to UD's definition <insert reference to UD website in sidenote>, they belong to \texttt{ADP}.\\

\paragraph{Tags with no UPOS counterparts}

% feature mappings
\subsection{Morphological Features}


% converting GRA codes to UD deprels
\subsection{Dependency Relations}

In CHILDES CHAT annotations, syntactic relations are represented using Grammatical Relations (GRs), first devised by \cite{Sagae2004}. Since GRs are dependency relations, they correspond to \texttt{deprel} in CoNLL-U annotations. However, CHILDES and UD uses not only different labels for the dependency relations, but also quite different parsing schemes such that sometimes a conversion from GR to deprel also introduces a change in the structure of the dependency tree.\\

Defining mappings between a GR and a UD deprel thus ranges from direct and one-to-one label translation to linguistic knowledge-based conditional decisions. This section provides explanations to the adaptations made and gives detailed examples for a selected few hard-to-decide cases.\\

Direct mappings are listed in the table below:
\begin{table}[h!]
\begin{tabular}{@{}lllll@{}}
\toprule
\textbf{CHILDES} & \textbf{UD} & \\ \midrule
SUBJ & nsubj & \\
CSUBJ & csubj & \\
OBJ & obj & \\
OBJ2 & iobj & \\
APP & appos & \\
AUX & aux & \\
CONJ & conj & \\
COORD & cc & \\
DET & det & \\
PUNCT & punct & \\
ROOT & root & \\\bottomrule
\end{tabular}
\end{table}


\subsection{PRED}
\begin{dependency}
	\begin{deptext}
	She \& is \& a \& student \& .\\
	PRON \& AUX \& DET \& NOUN \& PUNCT\\
	\end{deptext}
	\depedge{4}{1}{nsubj}
	\depedge{4}{2}{aux}
	\depedge{4}{3}{det}
	\depedge{4}{5}{punct}
	\deproot{4}{root}
\end{dependency}

\begin{dependency}
	\begin{deptext}
	She \& is \& a \& student \& .\\
	PRON \& AUX \& DET \& NOUN \& PUNCT\\
	\end{deptext}
	\depedge{2}{1}{SUBJ}
	\depedge{2}{4}{PRED}
	\depedge{4}{3}{DET}
	\depedge{2}{5}{PUNCT}
	\deproot{2}{ROOT}
\end{dependency}

\begin{dependency}
	\begin{deptext}
	She \& became \& happier \& .\\
	PRON \& VERB \& ADJ \& PUNCT\\
	\end{deptext}
	\depedge{2}{1}{nsubj}
	\depedge{2}{3}{xcomp}
	\depedge{2}{5}{punct}
	\deproot{2}{root}
\end{dependency}

\begin{dependency}
	\begin{deptext}
	She \& became \& happier \& .\\
	PRON \& VERB \& ADJ \& PUNCT\\
	\end{deptext}
	\depedge{2}{1}{SUBJ}
	\depedge{2}{3}{PRED}
	\depedge{2}{5}{PUNCT}
	\deproot{2}{ROOT}
\end{dependency}

\paragraph{\texttt{INCROOT}}

\paragraph{\texttt{PRED}}


\begin{table}[h!]
\begin{tabular}{@{}lllll@{}}
\toprule
\textbf{Universal POS Tags} & \textbf{Description}\\ \midrule
ADJ & adjective\\
ADV & adverb\\
INTJ & interjection\\
NOUN & noun\\
PROPN & proper noun\\
VERB & verb\\\midrule
ADP & adposition\\
AUX & auxiliary verb\\
CCONJ & coordinating conjunction\\
DET & determiner\\
NUM & numeral\\
PART & particle\\
PRON & pronoun\\
SCONJ & subordinating conjunction\\\midrule
PUNCT & punctuation\\
SYM & symbol\\
X & other\\\bottomrule
\end{tabular}
\end{table}

% \chapter{Ideas} % Main chapter title

\label{Chapter8} % For referencing the chapter elsewhere, use \ref{Chapter8}

\section{Morphosyntax in CHAT and CoNLL-U}

Both annotation formalisms support the storage of morphosyntactic information, although in different ways. In this chapter, I show the general structure of these formats using minimal examples, compare their ways of organising information and point out issues that I found worth discussing during the implementation process of chatconllu.

\subsection{Extracting Information from Dependent Tiers}

\paragraph{lemma} After looking at the syntax of the \%mor tier, one should realise that getting the lemma out of a MOR string is not an easy task. There are several pitfalls to avoid:\\
\begin{itemize}
	\item The semantics of the symbols may be ambiguous. Take the dash symbol as an example, for dashed words (words with dash symbols) like \texttt{tic-tac-toe}, it is merely a connector between parts of the word, but it could also appear as an indicator for suffixes, as in the example of \texttt{part|get-PRESP} for the word \emph{getting}.
	\item One should also notice that the stem is not the lemma, just that for English, these two share the same form in most cases. For instance, for \emph{untied}, we have the MOR segment \texttt{un\#v|tie-PASTP}, which means in plain English that the word has a stem \emph{tie} with prefix \emph{un}, its part-of-speech is verb and this word form is the past participle of the verb. On the other hand, in CoNLL-U, the lemma field is reserved for lemma, not word stem, therefore the prefix should be put back in place to produce the real lemma.
\end{itemize}


\texttt{prefix\#pos|stem\&fusionalsuffix-suffix=translation}

% \texttt{compound_pos|+component_pos|component_stem+component_pos|component_stem}

\texttt{prefix\#pos|stem\&fusionalsuffix-suffix=translation}


\section{Related Work}

\paragraph{pyconll} % (fold)
\label{par:pyconll}
\emph{pyconll} (\cite{pyconll}) is a low-level API for processing CoNLL-U formatted files. In this thesis, I used pyconll for parsing the converted conllu files and especially for reading the values of individual fields.

% \paragraph{CHAT2CONLLU and CONLL2CHAT} % (fold)
% \label{par:chat2conllu}
% \emph{CHAT2CONLL} and \emph{CONLL2CHAT} are officially part of the CLAN program that transforms CHAT files to CoNLL formats.

\paragraph{Dependency Parsing for Low-resource Spontaneous Speech} % (fold)
\label{par:zoey}
(\cite{liu2021}) used a low-resource setting, focusing on one child from the Brown Corpus (\cite{brown1973}) and tested the performance of dependency parsers on parent-child conversations. The out-domain parser was trained on written text and performed well on adult speech but not so well on child speech, while the parser trained on a limited amount of in-domain spoken data improved the parsing result of child speech to be comparable with adult speech. As part of their study, they created a semi-automatic adaptation of annotation formats from CHILDES to UD.

\paragraph{The AnnCor CHILDES Treebank}
(\cite{odijk2018anncor}) preprocessed the Dutch CHILDES corpora, augmented them with syntactic annotations using the Alpino parser, and performed partial manual verification. While creating this treebank, they noticed the many disadvantages of annotation discrepancies and argued that annotation guidelines should be explicit and should be strictly followed to produce consistency. They suggest combining human annotations with automatic checks by computer programs. They also pointed out that although the grammar used in child speech is limited since the child is still acquiring the language, assigning adult syntactic structures to their utterances is still valid.

%----------------------------------------------------------------------------------------
% THESIS CONTENT - APPENDICES
%----------------------------------------------------------------------------------------

% By using input instead of include for the chapters we are able to move the following line here
% Therefore the addition before the last chapter is not necessary anymore.

% Call the following chapters "Appendix" inside the table of contents
\addtocontents{toc}{\string\def\string\chaptername{Appendix}}

\appendix % Cue to tell LaTeX that the following "chapters" are Appendices

% Ensure proper section numbering in appendix, e.g., A.1, A.2, B.1, …
\renewcommand{\thesection}{\thechapter.\arabic{section}}
\renewcommand{\thesubsection}{\thesection.\arabic{subsection}}
\renewcommand{\thesubsubsection}{\thesubsection.\arabic{subsubsection}}

%%% CHANGES NEEDED HERE
%
% Include the appendices of the thesis as separate files from the Appendices folder
% Uncomment the lines as you write the Appendices

% Appendix A
 
\chapter{Morphosyntactic Mappings}
\label{appendixa}

\section{Part-of-Speech Mapping}

\begin{table}[h!]
\caption {Complete Set of Universal POS Tags} \label{tab:uposset}
\begin{tabularx}{\linewidth}{@{}llX@{}}
\toprule
 & \textbf{Universal POS Tag} & \textbf{Description}\\ \midrule
\textbf{Open class} & ADJ & adjective\\
\textbf{words} & ADV & adverb\\
 & INTJ & interjection\\
 & NOUN & noun\\
 & PROPN & proper noun\\
 & VERB & verb\\\addlinespace
\textbf{Closed class} & ADP & adposition\\
\textbf{words} & AUX & auxiliary verb\\
 & CCONJ & coordinating conjunction\\
 & DET & determiner\\
 & NUM & numeral\\
 & PART & particle\\
 & PRON & pronoun\\
 & SCONJ & subordinating conjunction\\\addlinespace
\textbf{Other} & PUNCT & punctuation\\
 & SYM & symbol\\
 & X & other\\\bottomrule
\end{tabularx}\\
\vspace{0.5cm}\\
All UPOS tags except \texttt{SYM} are used in chatconllu's POS mappings.\\
\end{table}

\newpage
\begin{table}[htp!]
\caption {chatconllu Part-of-Speech Mapping 1---open-class words} \label{tab:posmap1}
\centering
\begin{tabularx}{\linewidth}{@{}llX@{}}
\toprule
\textbf{Universal POS tag} & \textbf{MOR POS tag} & \textbf{Description}\\ \midrule
ADJ & adj & adjective\\
 & adj:pred & adjective-predicative\\\addlinespace
ADV & adv & adverb\\
 & adv:tem & adverb-temporal\\\addlinespace
INTJ & bab & babbling\\
 & co & communicator\\
 & fil & filler\\
 & sing & singing\\
 & wplay & wordplay\\\addlinespace
NOUN & n & noun\\
 & n:let & noun-letter\\
 & n:pt & noun-plurale tantum\\
 & on & onomatopoeia\\
 & onoma & onomatopoeia\\\addlinespace
PROPN & n:prop & proper noun\\\addlinespace
VERB & v & verb\\
 & vimp & verb imperative\\\bottomrule
\end{tabularx}\\
\vspace{0.5cm}
All MOR part-of-speech categories handled by the current version of chatconllu (part 1).\\
\end{table}

\newpage
\vspace{-2cm}
\begin{table}
\caption {chatconllu Part-of-Speech Mapping 2---closed-class words} \label{tab:posmap2}
\centering
\begin{tabularx}{\linewidth}{@{}llX@{}}
\toprule
\textbf{Universal POS tag} & \textbf{MOR POS tag} & \textbf{Description}\\ \midrule
ADP & post & postmodifier\\
 & prep & preposition\\
 & vpfx & preverb/verbal particles\\\addlinespace
AUX & aux & auxiliary\\
 & cop & verb-copula\\
 & mod & verb-modal\\
 & v:aux & verb-auxiliary\\\addlinespace
CCONJ & coord & coordinator\\\addlinespace
DET & art & article\\
 & det:art & determiner-article\\
 & det:dem & determiner-demonstrative\\
 & det:int & determiner-interrogative\\
 & det:num & determiner-numeral\\
 & det:poss & determiner-possessive\\
 & prepart & preposition with article\\
 & qn & quantifier\\
 & quant & quantifier\\\addlinespace
NUM & num & numeral\\\addlinespace
PART & inf & infinitive\\
 & neg & negative\\
 & part & particle\\\addlinespace
PRON & pro:dem & pronoun-demonstrative\\
 & pro:exist & pronoun-existential\\
 & pro:indef & pronoun-indefinitive\\
 & pro:int & pronoun-interrogative\\
 & pro:obj & pronoun-object\\
 & pro:per & pronoun-personal\\
 & pro:poss & pronoun-possessive\\
 & pro:refl & pronoun-reflexive\\
 & pro:rel & pronoun-relative\\
 & pro:sub & pronoun-subject\\\addlinespace
SCONJ & comp & complementizer\\
 & conj & conjunction\\\bottomrule
\end{tabularx}\\
\vspace{0.5cm}
All MOR part-of-speech categories handled by the current version of chatconllu (part 2).\\
\end{table}
\clearpage

\begin{table}[htp!]
\caption {chatconllu Part-of-Speech Mapping 3---other} \label{tab:posmap3}
\centering
\begin{tabularx}{\linewidth}{@{}llX@{}}
\toprule
\textbf{Universal POS tag} & \textbf{MOR POS tag} & \textbf{Description}\\ \midrule
PUNCT & beg & MOR punctuation mark \texttt{‡}\\
 & cm & MOR punctuation mark \texttt{,}\\
 & end & MOR punctuation mark \texttt{„}\\
 & bq & MOR punctuation mark \texttt{“}\\
 & eq & MOR punctuation mark \texttt{”}\\
 & bq2 & MOR punctuation mark \texttt{‘}\\
 & eq2 & MOR punctuation mark \texttt{’}\\\addlinespace
X & chi & child-invented form\\
 & dia & dialect\\
 & fam & family-specific form\\
 & meta & metalinguistic use\\
 & neo & neologism\\
 & none & none\\
 & phon & phonologically consistent form\\
 & test & test words like "wug"\\
 & uni & Unibet transcription\\
 & L2 & second-language form\\
\bottomrule
\end{tabularx}\\
\vspace{0.5cm}
All MOR part-of-speech categories handled by the current version of chatconllu (part 3).\\
\end{table}

\newpage
\section{Morphological Features Mapping}
\begin{table}[h!]
\caption {chatconllu Morpheme-Feature Mapping} \label{tab:featsmap}
\centering
\begin{tabularx}{\linewidth}{@{}lXl@{}}
\toprule
% \caption {chatconllu Morpheme-Feature Mapping} \label{tab:featsmap}
% \begin{tabular}{@{}lll@{}}
% \toprule
\textbf{Morpheme Code} & \textbf{UD feats}\\ \midrule
    cp & Degree=Cmp\\
    sp & Degree=Sup\\\addlinespace
    m & Gender=Masc\\
    f & Gender=Fem\\\addlinespace
    cond & Mood=Cnd\\
    imp & Mood=Imp\\
    sub & Mood=Sub\\
    subj & Mood=Sub\\\addlinespace
    pl & Number=Plur\\
    1p & Number=Plur|Person=1\\
    2p & Number=Plur|Person=2\\
    3p & Number=Plur|Person=3\\
    3sp & Number=Plur|Person=3\\
    sg & Number=Sing\\
    1s & Number=Sing|Person=1\\
    2s & Number=Sing|Person=2\\
    12s & Number=Sing|Person=2\\
    3s & Number=Sing|Person=3\\
    13s & Number=Sing|Person=3\\
    23s & Number=Sing|Person=3\\\addlinespace
    poss & Poss=Yes\\\addlinespace
    impf & Tense=Imp\\
    fut & Tense=Fut\\
    past & Tense=Past\\
    pret & Tense=Past\\
    pastp & Tense=Past|VerbForm=Part\\
    pp & Tense=Past|VerbForm=Part\\
    pres & Tense=Pres\\
    ppre & Tense=Pres|VerbForm=Part\\
    presp & Tense=Pres|VerbForm=Part\\\addlinespace
    inf & VerbForm=Inf\\\addlinespace
    pass & Voice=Pass\\\bottomrule
\end{tabularx}\\
\vspace{0.5cm}
All MOR grammatical morpheme codes handled by the current version of chatconllu.\\
\end{table}


\newpage
\section{Dependency Relation Mapping}\todo{change the table}
\begin{table}[h!]
\caption {chatconllu Dependency Relation Mapping} \label{tab:grmap}
\centering
\begin{tabularx}{\linewidth}{@{}lXl@{}}
\toprule
\textbf{GRA Grammatical Relation} & \textbf{UD deprel}\\ \midrule
    12s & Number=Sing|Person=2\\
    13s & Number=Sing|Person=3\\
    1p & Number=Plur|Person=1\\
    1s & Number=Sing|Person=1\\
    23s & Number=Sing|Person=3\\
    2p & Number=Plur|Person=2\\
    2s & Number=Sing|Person=2\\
    3p & Number=Plur|Person=3\\
    3s & Number=Sing|Person=3\\
    3sp & Number=Plur|Person=3\\
    cond & Mood=Cnd\\
    cp & Degree=Cmp\\
    f & Gender=Fem\\
    fem & Gender=Fem\\
    fut & Tense=Fut\\
    inf & VerbForm=Inf\\
    imp & Mood=Imp\\
    impf & Tense=Imp\\
    m & Gender=Masc\\
    pass & Voice=Pass\\
    past & Tense=Past\\
    pastp & Tense=Past|VerbForm=Part\\
    pl & Number=Plur\\
    poss & Poss=Yes\\
    pp & Tense=Past|VerbForm=Part\\
    ppre & Tense=Pres|VerbForm=Part\\
    pres & Tense=Pres\\
    presp & Tense=Pres|VerbForm=Part\\
    pret & Tense=Past\\
    sg & Number=Sing\\
    sp & Degree=Sup\\
    sub & Mood=Sub\\
    subj & Mood=Sub\\\bottomrule
\end{tabularx}\\
\vspace{0.5cm}
All GRASP GRs handled by the current version of chatconllu.\\
\end{table}

% \include{appendices/appendixB}
% \include{appendices/appendixC}



%----------------------------------------------------------------------------------------
% BIBLIOGRAPHY
%----------------------------------------------------------------------------------------

% Bibliography has no wide margins:
\newgeometry{
	inner=2cm, % Inner margin
	outer=2cm, % Outer margin
	marginparwidth=0cm,
	marginparsep=0mm,
	bindingoffset=.5cm, % Binding offset
	top=1.5cm, % Top margin
	bottom=2.5cm, % Bottom margin,
	includehead,
	includefoot
	% showframe, % Uncomment to show how the type block is set on the page
}

\addchap{References}

% enables two-column layout for bibliography
\setlength\columnsep{2em}
\begin{multicols}{2}
	\begin{refcontext}[sorting=nyt] % sort bibliography by last name, year, title
		\renewcommand*{\bibfont}{\small\RaggedRight}
		\linespread{1.0}\selectfont % increase linespread if desired (not recommended)
		\printbibliography[heading=none]
	\end{refcontext}
\end{multicols}

%----------------------------------------------------------------------------------------

%----------------------------------------------------------------------------------------
% DECLARATION PAGE
%----------------------------------------------------------------------------------------

\begin{declaration}
\addchaptertocentry{\authorshipname} % Add the declaration to the table of contents

% TODO Change the declaration according as needed. *

%\selectlanguage{ngerman}
Ich erkläre hiermit gemä\ss\ \S~17 Abs.\,2 APO, dass ich die vorstehende {\thesistype}arbeit selbständig\\ verfasst und keine anderen als die angegebenen Quellen und Hilfsmittel benutzt habe.

\bigskip
\bigskip

\begin{tabular}{@{}l@{}}
	T\"ubingen, den \rule[-0.8em]{10em}{0.5pt}\\[2ex]
	~
\end{tabular}
\hspace{\fill}%
\begin{tabular}{@{}c@{}}
	\rule[-0.8em]{20em}{0.5pt}\\[2ex]
	\authorname
\end{tabular}\hspace{\fill}

\end{declaration}

\end{document}
